\documentclass[12pt, a4paper]{article}
\usepackage[utf8]{inputenc}

\usepackage[a4paper,top=3cm,bottom=2cm,left=2.5cm,right=2.5cm,marginparwidth=2.55cm]{geometry}
%\usepackage[utf8]{inputenc}
%\usepackage[T1]{fontenc}
\usepackage{adjustbox}
\usepackage{amsmath}
\usepackage{amssymb}
\usepackage{authblk}
\usepackage[ngerman, english]{babel}
\usepackage{booktabs}
\usepackage{lineno}
\usepackage{csquotes}
\usepackage{float}
\usepackage{makecell}
\usepackage[maxfloats=256]{morefloats}
\usepackage{graphicx}
\usepackage{lmodern}
\usepackage{longtable}
\usepackage{lscape}
\usepackage[skip = 1pt, indent = 20pt]{parskip}
\usepackage{placeins}
\usepackage{pgffor}
\usepackage{caption}
\usepackage{subcaption}
\usepackage{setspace}
\usepackage{tabularx}
\usepackage{threeparttablex}
\usepackage[table]{xcolor}
\usepackage{xtab}

\usepackage{xr-hyper}


\usepackage[draft]{hyperref}

\externaldocument{main}

\hypersetup{
    %pdftitle={Heterogeneity in carbon intensity of consumption},
    % pdfpagemode=FullScreen,
    %linkbordercolor = 0 1 1
}

\maxdeadcycles = 1000

\onehalfspace

% everything concerning biblatex

\usepackage[backend = biber,
            style=authoryear-icomp,
            citestyle=authoryear-comp,
            url=false,
            isbn=false,
            doi=true,
            % doi=false,
            maxcitenames=2,
            uniquelist=false,
            date=year,
            uniquename=false,
            sorting=ynt,
            sortcites]{biblatex}

\addbibresource{5_Literature/Literature.bib} % TBD

% Setup
\graphicspath{{../1_Figures/}}

% new environment

\newenvironment{subcaption2}
{\strut
\vspace{-5pt}
\begin{minipage}[b]{0.95\textwidth}
  \hspace*{-\parindent}
  \footnotesize}
 {\end{minipage}}

% For continous counting in subfigures
\renewcommand*{\thesubfigure}{\thefigure.\arabic{subfigure}}

 % Title
\title{The heterogeneous effects of climate policy on households: Evidence from 88 countries\footnote{We thank Florian Nachtigall, Lennart Stern, Franziska Klein, Linus Mattauch, Michael Jakob, Karolina Rütten, Felix Knopp, and Thomas Sterner for helpful comments and suggestions. We thank Marlene Merchert, Charlotte Buder, and Paula Blechschmidt for great research assistance. Both authors acknowledge funding through the European Commission’s Horizon Europe Programme (101056873, ELEVATE and 101069880, AdJUST) and Deutsche Gesellschaft für Internationale Zusammenarbeit (GIZ, 812388773, Förderung sozial akzeptierter CO\textsubscript{2}-Bepreisung). Further, both authors acknowledge financial support used for data procurement from the Federal Ministry of Education and Research (BMBF, 03SFK5J0, Ariadne). Leonard Missbach acknowledges funding from the German Federal Environmental Foundation (DBU). Beyond, there are no further interested parties from which both authors have received significant financial support. Both authors, their close relatives or partners do not hold any paid or unpaid positions as officers, directors, or board members of relevant non-profit organizations or profit-making entities. No other party had a right to review and edit the manuscript prior to its circulation.}}

%\author[]{}
\author[1]{Leonard Missbach\footnote{Corresponding author: \href{mailto:leonard.missbach@pik-potsdam.de}{leonard.missbach@pik-potsdam.de}. Potsdam Institute for Climate Impact Research, P.O. Box 60 12 03, 14412 Potsdam, Germany}}
\author[1,2]{Jan Christoph Steckel}
\affil[1]{\small Potsdam Institute for Climate Impact Research (PIK), Potsdam, Germany}
\affil[2]{\small Technical University of Munich, Munich, Germany}

%\date{October 2024}

\begin{document}

%\maketitle

\vspace{30mm}

%\noindent A full version containing the image-intensive appendix can be found \href{https://www.econstor.eu/bitstream/10419/301069/1/Manuscript_20240815.pdf}{here}.


\thispagestyle{empty}

\linenumbers
\modulolinenumbers[5]

\setcounter{page}{1}

%TC:ignore
\appendix 
\begin{refsection}
{\Huge Online Appendix} \label{sec:appendix}

The heterogeneous effects of climate policy on households: Evidence from 88 countries

\clearpage
\section{Data cleaning} \label{sec:cleaning}

We describe our approach to collecting, cleaning, and harmonizing microdata, to feature engineering for machine learning analyses, and to the identification of descriptive country clusters.

\subsection{Collecting household data}\label{sec:cleaning_1}

We collect household budget survey data and extract several information before cleaning and harmonizing. Household budget survey data are often publicly available, but sometimes subject to a considerable fee. Table \ref{tab:datasets} provides publishing organizations, names of surveys, and links to datasets used in this study.

\begin{itemize}
    \item For each household, we include sociodemographic information about household members where available. In all survey, households are represented through 'household heads', i.e., persons that often contribute the largest share of household income or are responsible for purchase decisions. We use information on the 'household head' as a proxy for the entire household and collect information on education, gender, self-identified ethnicity, nationality, or religion of the 'household head'. We standardize information on education by using the International Standard Classification of Education (ISCED) to facilitate comparison across countries.
    \item We include spatial information where available, for example identifier for sub-national areas (provinces or states), sub-sub-national areas (districts) or villages. Often, surveys include an indicator for whether households live in urban or rural areas. Definitions of \textit{urban} and \textit{rural} may not be consistent across countries, but are certainly consistent within countries.
    \item We include information on energy use, such as on primary fuels used for cooking, lighting, and heating. We harmonize information on fuels across countries to account for different names and levels of detail across countries. For example, cooking fuels include charcoal, coal, electricity, firewood, gas, kerosene, liquid fuel, LPG, other biomass, or unknown fuels.
    \item We capture information on electricity access and create a binary variable that indicates if households have access to electricity through electricity grids, through generators, or solar panels. 
    \item We collect information about ownership of major transport vehicles (such as cars, motorcycles, and trucks) and major household appliances (such as refrigerators, air conditioners, washing machines, and television). For each country, we only include information about ownership, but not about the precise number of owned vehicles and appliances. This helps improve consistency across countries.
    \item We collect all available information on household-level expenditures, integrating information from household-level and individual-level diary entries. We do not include consumption information from home production, received as gifts, or as remuneration for labor. Our rationale is that it would be difficult for a climate policy instrument to cover self-produced goods and services that are not purchased on markets. We include all expenditures at the item level and extrapolate expenditures to yearly values. Often, households track expenditures over the course of a few weeks, but provide details on less frequent purchases in the past months or year. For more frequently purchased items, mostly food, this approach neglects seasonal consumption patterns, but resulting bias should be sufficiently small, since households are surveyed throughout the year.
    \item We do not include imputed expenditures, for example, for hypothetical rental payments, since including them would inaccurately overestimate total household expenditures and bias expenditure shares towards less carbon-intensive services.
\end{itemize}

Code written for each country-level dataset can be found in a stable online repository (see Appendix \ref{code}).

\subsection{Cleaning and harmonizing household data}\label{sec:cleaning_2}

Building on collected microdata from household budget surveys we perform several cleaning steps in order to harmonize datasets across countries as far as possible. Conditional on country-level data availability, we ensure that we clean all available data consistently. 

\begin{itemize}
    \item We remove households from the sample with missing information for key variables such as household size, sampling weights, or total expenditures.
    \item We code and treat missing information about other variables as missing or 'unknown' and remove variables for each country, if missing information are dominant.
    \item For each country, we address outliers of household expenditures at the item level. We consider any observation an outlier if it is in the 99\textsuperscript{th} percentile of all non-zero expenditures. We replace this observation with item-level median expenditures, thereby assuming that expenditure shares on such items are non-zero, but absolute values might have been exaggerated because of misreporting.
    \item We remove observations, if expenditures are negative, for example, because households sell items.
    \item We remove duplicates from our sample. We check separately for duplicates at the level of household information and at the level of expenditures on individual consumption items: We consider households spending the same amount of money on the same items duplicates.
    \item We remove all households from our sample if aggregate expenditures exceed mean aggregate expenditures by five standard deviations ($z>5$).
    \item We use inflation rates from \textcite{IMF.2020} and exchange rates from the \textcite{WorldBankGroup.2023} to convert all local currencies to USD for the year 2017. Expenditures from surveys conducted before 2017 are inflated; expenditures from surveys after 2017 are deflated. This ensures consistency with calculated sectoral CO\textsubscript{2} intensities as they refer to the year 2017. This approach does however neglect that expenditure shares may change with rising incomes and inflation. Adjustment for purchasing power parity (PPP) is not necessary because we refrain from comparing households across countries.
    \item We create matching tables to assign country-level expenditure items to 65 aggregate sectors and to four broad expenditure categories (energy, food, goods, services). Items, that are difficult to match to a specific sector or to a specific category, for example, 'other expenses', are matched to artificial sectors and categories labeled 'other'. Admittedly, we assume a carbon intensity of zero for such items in absence of more detailed information, but expenditure shares are generally low (0.7\% on average across country-level averages). 
    \item We delete observations for items indicating aggregate categories if this would lead to double-counting of single expenditures. We delete observations for items indicating taxes (e.g., 'property tax'), since including them would prove inaccurate to calculating expenditure shares and because items indicating tax payments are not available in each country.
    \item We match items pertaining to fuels such as firewood, charcoal, and other biomass to the sector \textit{lumber} to account for indirect emissions attributable to production, transportation, and retail of these goods. However, we treat direct CO\textsubscript{2} emissions of such fuels as zero, in line with assumptions by the IPCC \autocite{Grad.2023}, but also because direct emissions of such fuels are often difficult to regulate.
    \item We also identify items indicating energy use and create separate variables listing expenditures for different energy items such as electricity, gasoline, diesel, kerosene, LPG, natural gas, charcoal, hard coal, firewood, and other biomass. All matching tables are available through a separate stable online repository (see Appendix \ref{code}).
\end{itemize}

This procedure helps ensuring that sectoral expenditure shares are comparable across countries, even though not all surveys include information on the same number and detail of consumption items. We proceed with assigning households to expenditure quintiles based on total household expenditures per capita to account for differing expenditure shares in larger households. We use expenditure quintiles for our analyses in Figures 1 and \ref{fig:Quint}.

Tables \ref{tab:A1}, \ref{tab:A2}, \ref{tab:A3}, \ref{tab:A4_CF}, \ref{tab:A5_LF} and \ref{tab:A6} show summary statistics for our final harmonized dataset, grouped by country and by country and expenditure quintile in Tables \ref{tab:A2} and \ref{tab:A3}. 

% tracking and documenting removals

\subsection{Feature engineering} \label{sec:featureengineering}

Based on our harmonized dataset, we perform feature engineering on our variables (features) with the \texttt{R}-package \texttt{recipes} before performing analyses with BRT.

\begin{itemize}
    \item We exclude any feature with missing variation (for four countries).
    \item We exclude categorical feature with extremely high granularity (such as district-level identifiers) or no granularity (such as education, in some cases).
    \item We exclude any feature with missing values.
    \item We remove the minimum number of features necessary to avoid high levels of correlation ($r>0.9$) between all features.
    \item We code observations as 'other' for each feature (except province-level, district-level and urban/rural identifiers) that account for less than 5\% of all observations.
    \item All country-level feature sets include total household expenditures (in USD 2017) and household size. The minimum number of included features (including binary, categorical, and continuous features) is 4 (for Sweden) and the maximum number of included features is 17 (for Benin, Burkina Faso, Côte d'Ivoire, Guinea-Bissau, Senegal, and Togo).
\end{itemize}

\subsection{Identification of descriptive country clusters}\label{sec:country_clusters}

The relationship between household-level characteristics and carbon intensity of consumption is unique to each country, but also depends on the availability of granular data. We adjust for the importance of individual features by multiplying the importance of individual features and the country-level goodness of fit (R\textsuperscript{2}) to account for differences in available features across countries (see also Figure \subref{fig:data_visualisation_3}). This approach also helps to account for the aggregate performance of country-level models and allows for the better comparison of feature importance across countries. Our resulting measure of feature importance thus accounts for the contribution of individual features to explaining \textit{observed} values of carbon intensity rather than \textit{predicted} values of carbon intensity.

It is important to note that our approach of adjusting feature importance by overall predictive performance reduces the bias introduced by the limited availability of some features in the data at hand. Uncorrected feature importance values can be exaggerated when few features are available, so countries with few features may end up in the wrong cluster, mainly because the BRT model cannot help explain much of the variation in carbon intensity. Instead, our approach ensures that all \textit{observable} features contribute to clustering. Despite many structurally unobservable characteristics at the household level, our approach may be justified under the assumption that policy design can naturally focus only on observable household characteristics. This is particularly true for the design of compensation policies and when targeting errors are to be minimized.

Based on the (adjusted) feature importance for each country, we use the k-means algorithm for clustering. If features are missing in the data, we assume that their share of contribution is zero.

As a robustness check, we replace missing entries with the average values of all non-missing values. We use adjusted feature importance, i.e., we assume that the (imputed) contribution of unobserved features does not have a strong correlation with observed features.

We normalize all feature values to allow for comparison between features. If a feature is more (less) important in one country compared to all other countries, the processed feature value will be relatively high (low). If a feature is equally important in all countries, the processed feature values will be close to zero. We also include vertical distribution coefficients $\widehat{V_{r}}$ for clustering because our measure of feature importance of household expenditures does not capture the \textit{direction} of effects.

K-means clustering is an unsupervised machine learning method and helps to analyze clusters of observations that are most similar in many variables \textit{within} each cluster and least similar in many variables \textit{across} clusters \autocite{MacQueen.1967}. We examine the optimal number of clusters ($\{k \in \mathbb{N} \mid 2  \leq k \leq 20 \}$) using average silhouette widths \autocite{Rousseeuw.1987} for each cluster \textit{k}.

The silhouette width $s_{i}$ expresses the average Euclidean distance of each observation \textit{i} to all other observations within its cluster and for the average distance to observations from the nearest neighboring cluster. Silhouette widths closer to $1$ indicate a good fit of an observation to its cluster, and silhouette widths closer to $-1$ indicate a poor fit. The average silhouette width $\overline{s_{k}}$ for each cluster \textit{k} expresses how well all observations fit, on average, into each cluster.

Our approach yields $k = 10$ as the number of clusters that maximizes the average silhouette width. See Figures \subref{fig:G3_silhouette_2} and \subref{fig:G4_silhouette_2} for visualization. We also show the optimal number of clusters for k-means clustering based on unadjusted feature importance ($k = 6$) and on adjusted and imputed feature importance ($k=11$) here. See Figures \subref{fig:G1_silhouette}, \subref{fig:G2_silhouette}, \subref{fig:G1_silhouette_3}, and \subref{fig:G2_silhouette_3} for visualization. Using unadjusted feature importance for clustering changes the interpretation of clusters: Features of countries within the same cluster contribute similarly to explaining variation in carbon intensity, regardless of the availability of features in the data and the explanatory power of the model.

Within the same cluster, individual features are similarly important in predicting households' carbon intensity. For each cluster, we compute averages for each feature to allow us to examine differences between countries in different clusters.%We also investigate which country is most likely to all other countries in each cluster maximizing silhouette width within each cluster and use it to demonstrate cluster-level characteristics. 

Clusters differ in the importance of individual features. Figure \ref{fig:A_fig_3} shows the relative importance of different features across clusters, after ordering the clusters by cluster size. The largest cluster comprises 38 countries, where most features contribute relatively little to explaining variation in carbon intensity. These countries have in common that it is difficult to predict households' carbon intensity with the available data, after adjusting for for model performance. Despite low model performance and lacking information about energy, these results also point to highly idiosyncratic predictors of heterogeneous carbon intensity with important implications for policy design, as transfers differentiated based on observable characteristics (such as household expenditures) will not be effective in compensation the households most affected by climate policy.

% C mean carbon intensity and car ownership, D vertical differences, regressive, hh expenditures, B cooking fuel, E heating fuel, F progressive and vehicles,
This descriptive clustering approach helps identifying similarities and differences across countries. Countries in Cluster C stand out because of relatively high average carbon intensity and car ownership as an important feature; In Cluster D, variation in household expenditures is an comparably important feature and analyses exhibit comparably large vertical, regressive differences. In contrast, households consume less carbon-intensively in Cluster F and the consumption of richer households is more carbon-intensive than of poorer households. In general, features describing household energy use may be helpful predictors in some clusters (e.g., in Clusters B, C, E, and F), but not necessarily in all countries. For example, in countries in Cluster A such as \hyperref[fig:5b_BRA]{Brazil}, \hyperref[fig:5b_URY]{Argentina}, and \hyperref[fig:5b_ETH]{Ethiopia}, where the predictive power of the models is relatively high, the main energy fuels contribute relatively little to the prediction.

\subsection{Policy simulation}\label{sec:policysimulation}

We show that heterogeneity in the carbon intensity of consumption at the household level is equivalent to heterogeneity in household-level costs of climate policy, assuming that such instruments increase marginal costs of emitting CO\textsubscript{2} and that producers pass on such costs to consumers. This analysis thus disregards general-equilibrium-effects on both the supply- and the demand-side. In general, any climate policy instrument is conceivable in this exercise that leads to increasing costs in equivalence to embedded direct and indirect emissions, including (but not limited to) carbon pricing, fossil fuel subsidy removal, or subsidies for low-carbon fuels.

The carbon intensity of consumption $e_{i}$ consists of sectoral carbon intensities and household-level sectoral expenditure shares as shown in Equation 5. 

For example, consider the case of carbon pricing, which can be thought as a tax $\tau$ in USD/tCO\textsubscript{2}. The total absolute costs from carbon pricing equals direct and indirect carbon emissions embedded in household consumption $E_{i}$ multiplied with $\tau$. Computing total relative costs $c_{i}$ requires division by total household expenditures $C_{i}$:

\begin{equation}
    c_{i} = \frac{E_{i}*\tau}{C_{i}}
\end{equation}

Relative additional costs, i.e., the carbon pricing incidence $c_{i}$ can be expressed in \%, i.e., the fraction of absolute additional costs ($USD_{\tau}$) over total household expenditures ($USD_{i}$). $c_{i}$ is equivalent to our expression for carbon intensity of consumption $e_{i}$, scaled by a proportional factor $\tau$. If $e_{A}=2*e_{B}$, then $c_{A}=2*c_{B}$, assuming that $e_{A}$ and $e_{B}$ express the carbon intensity covering all nationally released CO\textsubscript{2} emissions for households \textit{A} and \textit{B} and that $c_{A}$ and $c_{B}$ refer to the relative carbon pricing incidence for a carbon price levied on all nationally released emissions in households \textit{A} and \textit{B}, respectively. In essence, heterogeneity in carbon intensity is equivalent to heterogeneity in household-level costs of climate policy instruments, under assumptions about how such instruments affect the marginal costs of emitting CO\textsubscript{2}.

In general, our modeling framework also allows for the simulation of other (sectoral) policies. Consider a carbon-tax-equivalent policy intervention in a specific sector, for example, in the transport sector, here denoted as $\tau_{s^{*}}$. Such a sector-specific tax would cover all direct and indirect emissions released in this sector $s^{*}$, but not emissions released in other sectors. Nevertheless, customer prices of goods and services from sectors other than transport would still increase because of embedded emissions from the transport sector.

Calculating additional sets of sectoral carbon intensities $e_{s^{*}}$ including direct and indirect emissions of different sectors can help simulate the impact of sectoral policies. Effectively, we only include direct and indirect CO\textsubscript{2} emissions released in sectors $s^{*}$.

It is also possible to analyze the distribution of regional policies, for example of carbon border taxes covering CO\textsubscript{2} emissions for imported goods and services. 

Supplementary Figure \ref{fig:comparison_policies} shows vertical and horizontal distribution coefficients for the national carbon intensity in all sectors, in the transport sector, in the electricity sector and for the international carbon intensity in all sectors. See also Table \ref{tab:A14}.

\clearpage

\renewcommand\thefigure{A.\arabic{figure}}
\renewcommand\thetable{A.\arabic{table}}
\setcounter{figure}{0}
\setcounter{table}{0}

\section{Supplementary figures} \label{sec:figures}

% B1: Overview of data work
\input{1_Figures/Figure_Datawork.tex}

% B2: Engel-curves
\input{1_Figures/Figure_Engel}
\clearpage

% B3: Carbon intensities
\input{1_Figures/Figure_Appendix_GTAP}
\clearpage

% B4: Boxplots over quintiles
\input{1_Figures/Figure_Appendix_Boxplots.tex}
\clearpage

% B5: Silhouette plots
\input{1_Figures/Figure_Silhouette}
\clearpage

% B6: Comparison of policies
\begin{figure}[ht!]
    \centering
    \includegraphics[width=\textwidth]{1_Figures/Figure 2/Figure_2_2017_Policy_C.pdf}
    %\includegraphics[width=\textwidth]{1_Figures/Figure 2/Figure_2_2017_Policy_B.pdf}
    \caption{Vertical and horizontal distribution coefficients for different policies}
    \label{fig:comparison_policies}
    \begin{subcaption2}
    This figure displays the vertical distribution coefficient comparing the median carbon intensity of the richest and the poorest quintile. The horizontal distribution coefficient compares the within-quintile differences (5\textsuperscript{th} to 95\textsuperscript{th} percentile within quintiles) of the richest and the poorest quintile. Rectangles (A) and (B) indicate higher carbon intensity (at the median) among the poorest quintile compared to the richest quintile; rectangles (C) and (D) indicate lower carbon intensity (at the median) among the poorest quintile compared to the richest quintile. Rectangles (A) and (C) indicate smaller within-quintile differences of carbon intensity among the poorest quintile compared to the richest quintile; rectangles (B) and (D) indicate larger within-quintile differences of carbon intensity among the poorest quintile compared to the richest quintile. Colors of points indicate GDP per capita for 2018 (in log-transformed constant 2010 USD).
    
    Panel 'National climate policy' shows the same values as Figure 2, i.e., distribution coefficients for carbon intensities accounting for all nationally released CO\textsubscript{2} emissions across all sectors. Panel 'International climate policy' shows distribution coefficients for carbon intensities accounting for global  CO\textsubscript{2} emissions embedded in national consumption. Panels 'Transport sector policy' and 'Electricity sector policy' display distribution coefficients for carbon intensities accounting for nationally released CO\textsubscript{2} emissions in the transport sector and electricity sector, respectively.
    \end{subcaption2}
\end{figure}
\clearpage

% B7: Comparison of R2
\begin{figure}[ht!]
    \centering
    \includegraphics[width=\textwidth]{1_Figures/Figures_Appendix/Figure_Comparison_Models.pdf}
    \caption{Goodness of fit (R\textsuperscript{2}) for rich and sparse BRT models}\label{fig:comparison}
    \label{fig:comparison_models}
    \begin{subcaption2}
    This figure shows goodness of fit (R\textsuperscript{2}) for sparse and rich boosted regression tree models. The sparse models include household expenditures as feature (blue point) and the rich models include all available features (red point), including household expenditures. We tune hyperparameters for each country and set of features and use fivefold cross-validation for evaluating model performance. See also Table \ref{tab:A8} for country-level MAE and RMSE.
    \end{subcaption2}
\end{figure}

\clearpage

% B8: unadjusted clustering
\input{1_Figures/Figure_4_Uncorrected}
\clearpage

%\setkeys{Gin}{draft}

% B9: SHAP-plots
\input{1_Figures/Figure_5b.tex}
\clearpage

%\setkeys{Gin}{draft=false}

% B10: Marginal effects-plots from Logit models
% \input{1_Figures/Figure_Marginal_Effects.tex}
% \clearpage

\begin{figure}[ht!]
    \centering
    \includegraphics{1_Figures/Figure 3/Figure_3_Corrected.pdf}
    \caption{Average feature importance across descriptive country clusters}
    \label{fig:A_fig_3}
    \begin{subcaption2}
    This figure shows the average importance of features (in normalized absolute average SHAP values) across all countries from each cluster A to J. Colors express the average importance of features in a cluster compared to other clusters. Blue (red) colors indicate that a feature is relatively less (more) important on average in a cluster compared to all other clusters. 'Sociodemographic' includes normalized absolute average SHAP values for features such as household size, education, gender, self-identified ethnicity, nationality, religion, and language. 'Spatial' includes normalized absolute SHAP values for province, district, and urban/rural identifiers. For vertical distribution, blue (red) colors indicate a lower (higher) median carbon intensity in the poorest quintile compared to the richest quintile. For average carbon intensity, blue (red) colors indicate a lower (higher) average carbon intensity across clusters.  For goodness of fit (R\textsuperscript{2}), blue (red) colors indicate a lower (higher) predictive performance compared to other countries. For the share of vertical heterogeneity in overall heterogeneity, blue (red) colors indicate a lower (higher) share across all countries. Average carbon intensity, the share of vertical heterogeneity in overall heterogeneity and R\textsuperscript{2} are not explicitly used for clustering. We assign countries to ten clusters using k-means clustering. We also show all values in Table \ref{tab:A9}.
    \end{subcaption2}
\end{figure}
\clearpage

%\setkeys{Gin}{draft}
% B11: Maps
%\input{1_Figures/Figure_Maps.tex}
% \clearpage

\begin{figure}[ht!]
    \centering
    \includegraphics[width=0.95\linewidth]{1_Figures/Figure 5b/Figure_5b_EU.jpg}
    \caption{Partial dependence plot (SHAP) for stylized analysis across 27 European countries}
    \label{fig:A_EU}
    \begin{subcaption2}
    This figure shows SHAP-values for predicting carbon intensity over feature values for 27 European countries. We have pooled country-level observations and included a country indicator in the feature set. This analysis can inform about the heterogeneous effects of a hypothetical European-level, cross-sectoral carbon price signal. We have coded missing feature values as explicitly missing. We have conducted hyperparameter tuning using fivefold cross-validation. We also use fivefold cross-validation for model evaluation. We calculate SHAP values for all observations, but show only a fraction of 20,000 (randomly selected) observations to improve legibility. The total number of observations is 268,616.

    The bar chart displays normalized average absolute SHAP-values for all features. Features with less than 3\% of normalized SHAP-values are subsumed as "Other features (Sum)". The right upper panel shows SHAP-values over total household expenditures over all countries. Colours represent household expenditures with blue (red) colors indicating lower (higher) household expenditures. The lower panel shows SHAP-values over household country. Colours represent household expenditures with blue(red) colors indicating lower (higher) household expenditures across all countries.

    Following this analysis, reducing overall heterogeneity would require country-differentiated transfers with households living in Poland, Czech Republic, Hungary, Slovakia, and Germany receiving relatively higher transfers. The right upper panel shows that relatively poorer households -- with household expenditures compared across countries -- would require relatively higher transfers than richer households.
    \end{subcaption2}
\end{figure}

\clearpage

%\setkeys{Gin}{draft=false}
\section{Supplementary tables} \label{sec:tables}

% Table C1: Household budget surveys
\input{2_Tables/Table_Datasets}
\clearpage

% Table C2: Summary statistics
\input{2_Tables/Table_Summary_A1}
\clearpage

% Table C3: Average expenditures and energy expenditure shares
\begingroup\fontsize{7}{9}\selectfont

\begin{ThreePartTable}
\begin{TableNotes}
\item \textit{Note: } 
\item This table shows average household expenditures and average energy expenditure shares for households in our sample. We estimate household-weighted averages for the whole population and per expenditure quintile.
\end{TableNotes}
\begin{longtable}[t]{l|rrrrrr|rrrrrrl|rrrrrr|rrrrrrl|rrrrrr|rrrrrrl|rrrrrr|rrrrrrl|rrrrrr|rrrrrrl|rrrrrr|rrrrrrl|rrrrrr|rrrrrrl|rrrrrr|rrrrrrl|rrrrrr|rrrrrrl|rrrrrr|rrrrrrl|rrrrrr|rrrrrrl|rrrrrr|rrrrrrl|rrrrrr|rrrrrr}
\caption{\label{tab:A2}Average household expenditures and average energy expenditure shares per expenditure quintile}\\
\toprule
\multicolumn{1}{c}{ } & \multicolumn{6}{c}{Average household expenditures [USD]} & \multicolumn{6}{c}{Average energy expenditure shares} \\
\cmidrule(l{3pt}r{3pt}){2-7} \cmidrule(l{3pt}r{3pt}){8-13}
\multicolumn{2}{c}{ } & \multicolumn{5}{c}{Expenditure quintile} & \multicolumn{1}{c}{ } & \multicolumn{5}{c}{Expenditure quintile} \\
\cmidrule(l{3pt}r{3pt}){3-7} \cmidrule(l{3pt}r{3pt}){9-13}
Country & All & EQ1 & EQ2 & EQ3 & EQ4 & EQ5 & All & EQ1 & EQ2 & EQ3 & EQ4 & EQ5\\
\midrule
\endfirsthead
\caption[]{Average household expenditures and average energy expenditure shares per expenditure quintile \textit{(continued)}}\\
\toprule
Country & All & EQ1 & EQ2 & EQ3 & EQ4 & EQ5 & All & EQ1 & EQ2 & EQ3 & EQ4 & EQ5\\
\midrule
\endhead

\endfoot
\bottomrule
\insertTableNotes
\endlastfoot
Argentina & 15,810 & 6,006 & 10,101 & 13,399 & 19,348 & 30,208 & 14\% & 17\% & 15\% & 14\% & 13\% & 10\%\\
Armenia & 4,779 & 1,788 & 2,698 & 3,410 & 4,486 & 11,516 & 19\% & 24\% & 21\% & 20\% & 18\% & 14\%\\
Australia & 64,969 & 29,372 & 44,957 & 59,601 & 75,879 & 115,068 & 7\% & 10\% & 8\% & 6\% & 5\% & 4\%\\
Austria & 38,472 & 22,665 & 30,220 & 35,378 & 41,890 & 62,212 & 10\% & 14\% & 11\% & 10\% & 9\% & 6\%\\
Bangladesh & 2,390 & 1,060 & 1,568 & 2,012 & 2,730 & 4,578 & 4\% & 4\% & 4\% & 4\% & 4\% & 4\%\\
Barbados & 17,650 & 7,207 & 12,753 & 16,956 & 19,867 & 31,427 & 13\% & 13\% & 13\% & 14\% & 13\% & 11\%\\
Belgium & 32,679 & 22,875 & 28,798 & 30,757 & 34,587 & 46,394 & 12\% & 15\% & 13\% & 12\% & 11\% & 9\%\\
Benin & 2,692 & 992 & 1,751 & 2,463 & 3,390 & 4,865 & 8\% & 6\% & 7\% & 8\% & 9\% & 11\%\\
Bolivia & 4,089 & 1,933 & 3,172 & 4,025 & 4,861 & 6,456 & 6\% & 7\% & 6\% & 6\% & 6\% & 6\%\\
Brazil & 10,916 & 2,581 & 5,127 & 7,755 & 11,913 & 27,207 & 14\% & 22\% & 15\% & 14\% & 12\% & 9\%\\
Bulgaria & 5,361 & 3,193 & 3,994 & 4,661 & 6,348 & 8,609 & 18\% & 20\% & 19\% & 19\% & 18\% & 15\%\\
Burkina Faso & 2,659 & 856 & 1,480 & 2,082 & 3,203 & 5,682 & 7\% & 4\% & 5\% & 6\% & 8\% & 11\%\\
Cambodia & 5,630 & 2,315 & 3,658 & 4,827 & 6,705 & 10,647 & 10\% & 12\% & 11\% & 10\% & 9\% & 9\%\\
Canada & 48,762 & 27,580 & 39,736 & 51,168 & 57,846 & 67,509 & 7\% & 9\% & 7\% & 7\% & 6\% & 5\%\\
Chile & 19,014 & 7,027 & 11,788 & 15,847 & 21,794 & 38,639 & 9\% & 13\% & 10\% & 9\% & 8\% & 6\%\\
Colombia & 6,855 & 1,573 & 3,031 & 4,479 & 7,130 & 18,063 & 9\% & 12\% & 10\% & 9\% & 7\% & 5\%\\
Costa Rica & 11,833 & 4,761 & 7,312 & 9,621 & 13,289 & 24,189 & 10\% & 13\% & 11\% & 10\% & 10\% & 8\%\\
Croatia & 11,967 & 7,622 & 9,658 & 11,423 & 13,626 & 17,519 & 18\% & 21\% & 20\% & 18\% & 17\% & 15\%\\
Cyprus & 26,609 & 14,972 & 21,888 & 26,210 & 31,629 & 38,365 & 14\% & 16\% & 15\% & 13\% & 12\% & 11\%\\
Czechia & 11,106 & 8,784 & 10,313 & 10,435 & 11,339 & 14,664 & 18\% & 20\% & 19\% & 19\% & 17\% & 15\%\\
Côte d’Ivoire & 3,241 & 1,426 & 2,384 & 3,220 & 3,981 & 5,194 & 6\% & 5\% & 6\% & 6\% & 6\% & 7\%\\
Denmark & 37,953 & 30,438 & 35,459 & 34,052 & 39,350 & 50,488 & 12\% & 13\% & 12\% & 12\% & 11\% & 9\%\\
Dominican Republic & 7,547 & 4,027 & 5,718 & 6,938 & 8,310 & 12,741 & 10\% & 9\% & 9\% & 9\% & 9\% & 12\%\\
Ecuador & 6,831 & 2,598 & 4,384 & 5,672 & 7,473 & 14,031 & 7\% & 8\% & 6\% & 6\% & 6\% & 7\%\\
Egypt & 2,449 & 1,818 & 2,254 & 2,503 & 2,679 & 2,992 & 6\% & 6\% & 6\% & 6\% & 6\% & 7\%\\
El Salvador & 5,757 & 1,288 & 2,977 & 4,740 & 6,944 & 12,835 & 20\% & 26\% & 23\% & 20\% & 17\% & 14\%\\
Estonia & 12,003 & 5,506 & 8,105 & 10,768 & 13,579 & 22,068 & 15\% & 19\% & 17\% & 15\% & 14\% & 11\%\\
Ethiopia & 1,168 & 315 & 637 & 894 & 1,508 & 2,485 & 5\% & 3\% & 4\% & 4\% & 7\% & 7\%\\
Finland & 31,858 & 22,862 & 26,643 & 29,442 & 33,002 & 47,374 & 8\% & 10\% & 9\% & 8\% & 7\% & 6\%\\
France & 29,109 & 17,162 & 24,209 & 28,678 & 31,849 & 43,643 & 10\% & 13\% & 11\% & 11\% & 9\% & 7\%\\
Georgia & 2,436 & 1,200 & 1,877 & 2,259 & 2,804 & 4,038 & 15\% & 16\% & 16\% & 16\% & 15\% & 14\%\\
Germany & 29,536 & 25,612 & 28,749 & 28,650 & 29,985 & 34,685 & 13\% & 15\% & 14\% & 13\% & 12\% & 10\%\\
Ghana & 2,381 & 1,152 & 1,939 & 2,414 & 2,942 & 3,457 & 8\% & 6\% & 8\% & 8\% & 9\% & 9\%\\
Greece & 19,197 & 11,054 & 14,310 & 17,371 & 20,712 & 32,546 & 14\% & 17\% & 16\% & 14\% & 12\% & 10\%\\
Guatemala & 5,678 & 2,573 & 3,998 & 5,080 & 6,481 & 10,266 & 16\% & 20\% & 16\% & 15\% & 15\% & 14\%\\
Guinea-Bissau & 3,690 & 1,509 & 2,510 & 3,345 & 4,413 & 6,678 & 4\% & 2\% & 2\% & 3\% & 5\% & 8\%\\
Hungary & 8,367 & 5,502 & 7,132 & 8,025 & 9,386 & 11,795 & 20\% & 22\% & 21\% & 21\% & 20\% & 17\%\\
India & 2,735 & 1,770 & 2,261 & 2,628 & 3,084 & 3,932 & 12\% & 11\% & 12\% & 12\% & 13\% & 12\%\\
Indonesia & 2,837 & 1,098 & 1,813 & 2,483 & 3,404 & 5,388 & 12\% & 14\% & 12\% & 12\% & 11\% & 11\%\\
Iraq & 14,016 & 5,818 & 9,100 & 11,805 & 15,711 & 27,645 & 9\% & 12\% & 10\% & 9\% & 8\% & 6\%\\
Ireland & 33,761 & 20,904 & 28,300 & 33,702 & 39,238 & 46,665 & 13\% & 16\% & 15\% & 13\% & 13\% & 10\%\\
Israel & 39,038 & 19,943 & 29,933 & 37,968 & 46,092 & 61,269 & 8\% & 10\% & 8\% & 7\% & 7\% & 5\%\\
Italy & 23,953 & 12,917 & 19,068 & 23,315 & 28,096 & 36,371 & 14\% & 19\% & 16\% & 14\% & 13\% & 10\%\\
Jordan & 11,929 & 7,223 & 9,466 & 11,178 & 13,911 & 17,879 & 18\% & 15\% & 16\% & 18\% & 19\% & 20\%\\
Kenya & 2,468 & 680 & 1,392 & 2,090 & 2,914 & 5,264 & 6\% & 6\% & 6\% & 7\% & 6\% & 6\%\\
Latvia & 10,314 & 5,205 & 6,954 & 8,672 & 11,423 & 19,331 & 17\% & 19\% & 20\% & 17\% & 17\% & 14\%\\
Liberia & 2,567 & 877 & 1,691 & 2,488 & 3,409 & 4,372 & 3\% & 3\% & 2\% & 3\% & 4\% & 5\%\\
Lithuania & 8,890 & 5,290 & 6,516 & 7,709 & 10,560 & 14,381 & 18\% & 18\% & 18\% & 19\% & 19\% & 16\%\\
Luxembourg & 49,872 & 32,803 & 40,602 & 50,226 & 56,769 & 68,992 & 9\% & 12\% & 10\% & 8\% & 8\% & 6\%\\
Malawi & 707 & 165 & 358 & 531 & 812 & 1,671 & 2\% & 0\% & 1\% & 2\% & 4\% & 6\%\\
Maldives & 20,194 & 10,575 & 15,911 & 19,808 & 24,853 & 29,857 & 7\% & 10\% & 8\% & 7\% & 5\% & 4\%\\
Mali & 3,452 & 1,195 & 2,032 & 2,987 & 4,420 & 6,630 & 6\% & 4\% & 6\% & 6\% & 8\% & 8\%\\
Mexico & 5,928 & 2,352 & 3,954 & 5,158 & 6,696 & 11,481 & 11\% & 9\% & 10\% & 11\% & 12\% & 11\%\\
Mongolia & 5,673 & 3,238 & 4,547 & 5,181 & 6,075 & 9,325 & 11\% & 12\% & 12\% & 12\% & 10\% & 8\%\\
Morocco & 7,397 & 3,925 & 5,379 & 6,478 & 8,183 & 13,020 & 8\% & 10\% & 8\% & 7\% & 7\% & 7\%\\
Mozambique & 2,743 & 247 & 789 & 1,610 & 3,362 & 7,706 & 3\% & 1\% & 1\% & 2\% & 4\% & 8\%\\
Myanmar (Burma) & 2,246 & 1,031 & 1,523 & 1,988 & 2,609 & 4,084 & 5\% & 5\% & 5\% & 5\% & 6\% & 6\%\\
Netherlands & 34,677 & 28,492 & 32,145 & 31,995 & 34,933 & 45,821 & 10\% & 12\% & 11\% & 10\% & 9\% & 8\%\\
Nicaragua & 4,816 & 1,410 & 2,557 & 3,608 & 5,262 & 11,248 & 6\% & 4\% & 5\% & 6\% & 7\% & 8\%\\
Niger & 1,900 & 620 & 1,108 & 1,504 & 2,106 & 4,162 & 3\% & 1\% & 1\% & 2\% & 3\% & 7\%\\
Nigeria & 2,311 & 1,064 & 1,788 & 2,326 & 2,933 & 3,444 & 5\% & 4\% & 4\% & 5\% & 6\% & 6\%\\
Norway & 53,153 & 28,948 & 41,897 & 51,022 & 60,273 & 83,666 & 10\% & 14\% & 12\% & 10\% & 9\% & 7\%\\
Pakistan & 3,512 & 2,121 & 2,732 & 3,124 & 3,829 & 5,756 & 9\% & 7\% & 8\% & 10\% & 10\% & 11\%\\
Paraguay & 7,393 & 2,467 & 4,802 & 6,951 & 9,082 & 13,665 & 10\% & 10\% & 11\% & 10\% & 11\% & 10\%\\
Peru & 4,675 & 1,603 & 3,123 & 4,354 & 5,618 & 8,678 & 8\% & 9\% & 9\% & 8\% & 8\% & 7\%\\
Philippines & 4,467 & 1,797 & 2,725 & 3,825 & 5,346 & 8,643 & 6\% & 4\% & 5\% & 6\% & 7\% & 7\%\\
Poland & 12,777 & 7,039 & 9,020 & 10,491 & 13,270 & 24,064 & 15\% & 16\% & 17\% & 16\% & 15\% & 10\%\\
Portugal & 17,659 & 8,994 & 12,929 & 16,248 & 20,121 & 30,016 & 17\% & 22\% & 19\% & 17\% & 15\% & 12\%\\
Romania & 9,698 & 3,724 & 5,731 & 8,802 & 12,591 & 17,646 & 14\% & 17\% & 17\% & 15\% & 12\% & 9\%\\
Russia & 6,293 & 2,948 & 4,511 & 5,557 & 6,822 & 11,632 & 2\% & 2\% & 2\% & 2\% & 2\% & 2\%\\
Rwanda & 1,261 & 409 & 674 & 921 & 1,368 & 2,933 & 3\% & 1\% & 2\% & 3\% & 4\% & 6\%\\
Senegal & 6,702 & 3,067 & 5,043 & 6,839 & 8,204 & 10,359 & 5\% & 3\% & 4\% & 5\% & 6\% & 7\%\\
Serbia & 7,608 & 4,582 & 6,653 & 7,751 & 8,868 & 10,187 & 14\% & 13\% & 15\% & 14\% & 15\% & 14\%\\
Slovakia & 12,838 & 8,803 & 10,974 & 11,969 & 13,512 & 18,939 & 20\% & 23\% & 21\% & 21\% & 18\% & 14\%\\
South Africa & 6,889 & 1,742 & 2,841 & 3,934 & 6,643 & 19,287 & 11\% & 11\% & 10\% & 11\% & 12\% & 12\%\\
Spain & 23,841 & 12,242 & 18,700 & 23,988 & 28,615 & 35,663 & 12\% & 14\% & 13\% & 12\% & 11\% & 9\%\\
Suriname & 7,588 & 2,945 & 5,058 & 6,844 & 8,982 & 14,125 & 6\% & 8\% & 7\% & 6\% & 5\% & 4\%\\
Sweden & 30,537 & 21,280 & 26,982 & 28,946 & 32,280 & 43,232 & 11\% & 13\% & 12\% & 11\% & 9\% & 8\%\\
Switzerland & 76,253 & 59,429 & 68,887 & 74,039 & 80,478 & 98,444 & 4\% & 5\% & 4\% & 4\% & 4\% & 3\%\\
Taiwan & 20,688 & 13,197 & 17,887 & 20,626 & 23,591 & 28,142 & 10\% & 11\% & 11\% & 11\% & 10\% & 9\%\\
Thailand & 3,749 & 1,037 & 1,873 & 2,999 & 4,748 & 8,088 & 20\% & 20\% & 23\% & 23\% & 19\% & 14\%\\
Togo & 2,382 & 818 & 1,539 & 2,282 & 3,154 & 4,117 & 8\% & 4\% & 7\% & 8\% & 9\% & 10\%\\
Turkey & 9,991 & 4,955 & 6,968 & 8,976 & 11,139 & 17,917 & 11\% & 11\% & 12\% & 12\% & 12\% & 10\%\\
Uganda & 1,262 & 288 & 656 & 1,036 & 1,606 & 2,724 & 5\% & 4\% & 3\% & 5\% & 6\% & 7\%\\
United Kingdom & 35,298 & 15,959 & 25,044 & 32,860 & 40,579 & 62,072 & 10\% & 12\% & 12\% & 10\% & 9\% & 7\%\\
United States & 43,763 & 24,302 & 34,000 & 41,714 & 48,683 & 70,156 & 10\% & 13\% & 12\% & 10\% & 9\% & 6\%\\
Uruguay & 21,055 & 8,143 & 13,359 & 18,383 & 24,905 & 40,498 & 10\% & 13\% & 11\% & 9\% & 8\% & 7\%\\
Vietnam & 2,362 & 762 & 1,435 & 2,072 & 2,975 & 4,566 & 5\% & 5\% & 5\% & 5\% & 5\% & 4\%\\*
\end{longtable}
\end{ThreePartTable}
\endgroup{}

\clearpage

% Table C4: Cooking fuel use
\input{2_Tables/Table_Summary_A4_Cooking_Fuels}
\clearpage

% Table C5: Lighting fuel use
\input{2_Tables/Table_Summary_A4_Lighting_Fuels}
\clearpage

% Table C6: Appliance ownership
\input{2_Tables/Table_Summary_A5_Appliances}
\clearpage

% \foreach \country in {ARG, ARM, BEL, BEN, BFA,
% BGD, BGR, BOL, BRA, BRB, CHL, %CIV,
%  COL, CRI, CYP, CZE, DEU, DNK, DOM, ECU, ESP, EST, ETH, FIN, FRA, GHA, GRC, GTM,
%  HRV, HUN, IDN, IND, IRL, IRQ, ISR, ITA, KEN, KHM, LBR, LTU, LUX, LVA, MAR, MDV, MEX, MLI, MMR, MNG, MWI, NER, NGA,  NIC, NLD, NOR, PAK, PER, PHL, POL, PRT, PRY, ROU, RWA, SEN, SLV, SUR, SVK, SWE, TGO, THA, TUR, UGA, URY, ZAF}{
%   \input{../2_Tables/3c_OLS_Logit_combined/Table_CF_CI_UL20_\country}
% }

% \clearpage

% Table C7: Carbon footprints and carbon intensities
\begingroup\fontsize{9}{11}\selectfont

\begin{ThreePartTable}
\begin{TableNotes}
\item \textit{Note: } 
\item This table shows average carbon footprints in tCO$_{2}$ and average carbon intensity in kgCO$_{2}$/USD for households in all countries of our sample. We estimate household-weighted averages for the whole population and per expenditure quintile.
\end{TableNotes}
\begin{longtable}[t]{l|rrrrrr|rrrrrrl|rrrrrr|rrrrrrl|rrrrrr|rrrrrrl|rrrrrr|rrrrrrl|rrrrrr|rrrrrrl|rrrrrr|rrrrrrl|rrrrrr|rrrrrrl|rrrrrr|rrrrrrl|rrrrrr|rrrrrrl|rrrrrr|rrrrrrl|rrrrrr|rrrrrrl|rrrrrr|rrrrrrl|rrrrrr|rrrrrr}
\caption{\label{tab:A3}Average carbon footprint and average carbon intensity per expenditure quintile}\\
\toprule
\multicolumn{1}{c}{ } & \multicolumn{6}{c}{Average carbon footprint [tCO$_{2}$]} & \multicolumn{6}{c}{Average carbon intensity [kgCO$_{2}$/USD]} \\
\cmidrule(l{3pt}r{3pt}){2-7} \cmidrule(l{3pt}r{3pt}){8-13}
\multicolumn{2}{c}{ } & \multicolumn{5}{c}{Expenditure quintile} & \multicolumn{1}{c}{ } & \multicolumn{5}{c}{Expenditure quintile} \\
\cmidrule(l{3pt}r{3pt}){3-7} \cmidrule(l{3pt}r{3pt}){9-13}
Country & All & EQ1 & EQ2 & EQ3 & EQ4 & EQ5 & All & EQ1 & EQ2 & EQ3 & EQ4 & EQ5\\
\midrule
\endfirsthead
\caption[]{Average carbon footprint and average carbon intensity per expenditure quintile \textit{(continued)}}\\
\toprule
Country & All & EQ1 & EQ2 & EQ3 & EQ4 & EQ5 & All & EQ1 & EQ2 & EQ3 & EQ4 & EQ5\\
\midrule
\endhead

\endfoot
\bottomrule
\insertTableNotes
\endlastfoot
Argentina & 16.4 & 9.6 & 13.5 & 15.9 & 19.7 & 23.5 & 1.27 & 1.78 & 1.42 & 1.24 & 1.07 & 0.82\\
Armenia & 3.4 & 1.7 & 2.5 & 3.2 & 3.7 & 5.8 & 0.87 & 0.99 & 0.95 & 0.95 & 0.85 & 0.62\\
Australia & 32.3 & 20.4 & 26.4 & 32.0 & 38.1 & 44.8 & 0.58 & 0.74 & 0.62 & 0.56 & 0.53 & 0.42\\
Austria & 10.5 & 8.1 & 9.6 & 10.1 & 11.4 & 13.4 & 0.31 & 0.37 & 0.33 & 0.30 & 0.29 & 0.24\\
Bangladesh & 0.8 & 0.3 & 0.5 & 0.6 & 0.9 & 1.6 & 0.33 & 0.34 & 0.31 & 0.31 & 0.33 & 0.35\\
Barbados & 9.8 & 4.2 & 7.2 & 10.8 & 12.4 & 14.5 & 0.58 & 0.56 & 0.59 & 0.63 & 0.60 & 0.50\\
Belgium & 13.4 & 10.9 & 12.8 & 13.3 & 14.0 & 16.1 & 0.44 & 0.49 & 0.47 & 0.46 & 0.43 & 0.37\\
Benin & 1.2 & 0.4 & 0.7 & 1.0 & 1.2 & 2.9 & 0.40 & 0.37 & 0.38 & 0.37 & 0.36 & 0.51\\
Bolivia & 1.5 & 0.7 & 1.3 & 1.6 & 1.9 & 2.2 & 0.39 & 0.41 & 0.40 & 0.40 & 0.40 & 0.36\\
Brazil & 4.6 & 1.4 & 2.5 & 3.6 & 5.3 & 10.3 & 0.48 & 0.60 & 0.49 & 0.47 & 0.44 & 0.39\\
Bulgaria & 3.1 & 1.9 & 2.3 & 2.9 & 3.8 & 4.6 & 0.57 & 0.56 & 0.56 & 0.60 & 0.59 & 0.55\\
Burkina Faso & 1.3 & 0.3 & 0.6 & 0.8 & 1.3 & 3.4 & 0.39 & 0.30 & 0.35 & 0.36 & 0.40 & 0.54\\
Cambodia & 2.6 & 1.2 & 1.8 & 2.2 & 3.0 & 4.6 & 0.47 & 0.48 & 0.51 & 0.46 & 0.45 & 0.44\\
Canada & 25.4 & 17.1 & 21.6 & 26.9 & 29.9 & 31.3 & 0.53 & 0.62 & 0.53 & 0.53 & 0.52 & 0.46\\
Chile & 8.3 & 4.2 & 6.2 & 7.8 & 9.6 & 13.8 & 0.50 & 0.63 & 0.54 & 0.50 & 0.44 & 0.37\\
Colombia & 2.8 & 1.0 & 1.7 & 2.2 & 3.1 & 6.1 & 0.50 & 0.62 & 0.56 & 0.51 & 0.44 & 0.35\\
Costa Rica & 2.6 & 1.0 & 1.8 & 2.3 & 3.2 & 4.6 & 0.22 & 0.21 & 0.24 & 0.23 & 0.23 & 0.20\\
Croatia & 5.0 & 2.9 & 4.3 & 4.9 & 5.8 & 7.1 & 0.40 & 0.34 & 0.41 & 0.41 & 0.41 & 0.43\\
Cyprus & 6.6 & 4.4 & 6.1 & 6.5 & 7.6 & 8.3 & 0.27 & 0.30 & 0.29 & 0.27 & 0.26 & 0.24\\
Czechia & 11.7 & 9.9 & 11.1 & 11.5 & 11.7 & 14.4 & 1.11 & 1.18 & 1.14 & 1.16 & 1.07 & 1.00\\
Côte d’Ivoire & 0.8 & 0.3 & 0.5 & 0.7 & 0.9 & 1.7 & 0.22 & 0.18 & 0.21 & 0.21 & 0.21 & 0.28\\
Denmark & 6.7 & 6.0 & 6.5 & 6.0 & 6.9 & 8.0 & 0.18 & 0.20 & 0.18 & 0.17 & 0.17 & 0.16\\
Dominican Republic & 2.6 & 1.1 & 1.6 & 2.1 & 2.6 & 5.4 & 0.30 & 0.27 & 0.28 & 0.30 & 0.30 & 0.38\\
Ecuador & 2.7 & 1.0 & 1.5 & 2.0 & 3.0 & 6.0 & 0.40 & 0.45 & 0.36 & 0.35 & 0.39 & 0.43\\
Egypt & 2.3 & 1.6 & 2.0 & 2.2 & 2.5 & 3.3 & 0.94 & 0.93 & 0.89 & 0.90 & 0.94 & 1.05\\
El Salvador & 1.9 & 0.7 & 1.3 & 1.7 & 2.1 & 3.7 & 0.37 & 0.51 & 0.42 & 0.34 & 0.30 & 0.27\\
Estonia & 5.8 & 3.4 & 4.6 & 5.7 & 6.3 & 8.9 & 0.54 & 0.63 & 0.60 & 0.54 & 0.49 & 0.43\\
Ethiopia & 0.1 & 0.0 & 0.1 & 0.1 & 0.1 & 0.2 & 0.09 & 0.12 & 0.10 & 0.09 & 0.08 & 0.08\\
Finland & 10.2 & 7.8 & 8.9 & 10.2 & 10.7 & 13.1 & 0.33 & 0.35 & 0.33 & 0.34 & 0.32 & 0.28\\
France & 8.9 & 6.3 & 8.2 & 9.2 & 9.2 & 11.4 & 0.32 & 0.39 & 0.35 & 0.33 & 0.29 & 0.26\\
Georgia & 2.5 & 1.2 & 2.0 & 2.5 & 3.0 & 3.8 & 0.96 & 0.92 & 0.98 & 1.00 & 1.00 & 0.93\\
Germany & 15.4 & 15.0 & 15.9 & 15.2 & 15.0 & 15.5 & 0.53 & 0.60 & 0.57 & 0.54 & 0.51 & 0.45\\
Ghana & 0.6 & 0.2 & 0.3 & 0.5 & 0.8 & 1.3 & 0.20 & 0.13 & 0.16 & 0.19 & 0.24 & 0.29\\
Greece & 7.2 & 4.8 & 6.0 & 7.0 & 7.8 & 10.5 & 0.40 & 0.44 & 0.43 & 0.41 & 0.38 & 0.34\\
Guatemala & 2.1 & 0.3 & 0.8 & 1.5 & 2.5 & 5.2 & 0.28 & 0.12 & 0.18 & 0.28 & 0.37 & 0.47\\
Guinea-Bissau & 0.9 & 0.3 & 0.4 & 0.6 & 0.9 & 2.2 & 0.19 & 0.15 & 0.15 & 0.17 & 0.18 & 0.27\\
Hungary & 5.7 & 3.6 & 5.0 & 5.8 & 6.7 & 7.6 & 0.70 & 0.65 & 0.72 & 0.74 & 0.72 & 0.66\\
India & 3.6 & 2.0 & 2.8 & 3.4 & 4.2 & 5.6 & 1.28 & 1.13 & 1.21 & 1.28 & 1.35 & 1.39\\
Indonesia & 4.0 & 1.6 & 2.6 & 3.6 & 4.9 & 7.5 & 1.45 & 1.45 & 1.46 & 1.45 & 1.46 & 1.44\\
Iraq & 11.5 & 6.4 & 9.2 & 11.1 & 13.2 & 17.9 & 0.97 & 1.19 & 1.05 & 0.99 & 0.88 & 0.72\\
Ireland & 13.6 & 10.0 & 12.7 & 13.8 & 15.7 & 15.7 & 0.45 & 0.53 & 0.49 & 0.44 & 0.43 & 0.36\\
Israel & 13.3 & 9.7 & 12.2 & 13.5 & 14.9 & 16.3 & 0.39 & 0.55 & 0.42 & 0.36 & 0.34 & 0.27\\
Italy & 9.2 & 5.8 & 8.0 & 9.2 & 10.6 & 12.5 & 0.41 & 0.46 & 0.43 & 0.40 & 0.39 & 0.36\\
Jordan & 13.7 & 6.6 & 9.9 & 12.8 & 16.7 & 22.5 & 1.10 & 0.93 & 1.03 & 1.14 & 1.19 & 1.24\\
Kenya & 1.1 & 0.2 & 0.5 & 0.8 & 1.3 & 2.9 & 0.40 & 0.37 & 0.36 & 0.38 & 0.42 & 0.49\\
Latvia & 4.7 & 3.4 & 3.7 & 3.9 & 5.4 & 7.0 & 0.52 & 0.73 & 0.56 & 0.45 & 0.49 & 0.38\\
Liberia & 0.6 & 0.1 & 0.3 & 0.6 & 0.9 & 1.4 & 0.19 & 0.11 & 0.15 & 0.20 & 0.22 & 0.27\\
Lithuania & 3.9 & 2.1 & 2.7 & 3.3 & 5.0 & 6.4 & 0.39 & 0.35 & 0.35 & 0.38 & 0.43 & 0.42\\
Luxembourg & 14.0 & 12.5 & 13.2 & 13.9 & 15.0 & 15.5 & 0.32 & 0.42 & 0.36 & 0.30 & 0.30 & 0.24\\
Malawi & 0.0 & 0.0 & 0.0 & 0.0 & 0.0 & 0.2 & 0.03 & 0.01 & 0.02 & 0.02 & 0.03 & 0.07\\
Maldives & 4.8 & 3.0 & 4.3 & 4.9 & 5.5 & 6.0 & 0.26 & 0.30 & 0.28 & 0.26 & 0.23 & 0.20\\
Mali & 1.1 & 0.4 & 0.6 & 0.8 & 1.3 & 2.2 & 0.27 & 0.25 & 0.27 & 0.25 & 0.28 & 0.31\\
Mexico & 5.5 & 2.1 & 3.7 & 4.9 & 6.6 & 10.3 & 0.93 & 0.87 & 0.93 & 0.96 & 0.98 & 0.92\\
Mongolia & 4.8 & 3.5 & 4.5 & 4.9 & 5.1 & 5.8 & 0.98 & 1.13 & 1.04 & 1.00 & 0.93 & 0.77\\
Morocco & 4.6 & 2.6 & 3.3 & 3.9 & 4.9 & 8.1 & 0.62 & 0.68 & 0.62 & 0.60 & 0.59 & 0.60\\
Mozambique & 1.0 & 0.0 & 0.1 & 0.3 & 0.7 & 3.8 & 0.22 & 0.23 & 0.14 & 0.17 & 0.20 & 0.36\\
Myanmar (Burma) & 1.2 & 0.4 & 0.6 & 0.9 & 1.3 & 2.7 & 0.45 & 0.36 & 0.40 & 0.41 & 0.46 & 0.62\\
Netherlands & 12.1 & 11.3 & 11.8 & 11.2 & 11.6 & 14.4 & 0.36 & 0.41 & 0.38 & 0.36 & 0.33 & 0.32\\
Nicaragua & 1.2 & 0.1 & 0.3 & 0.7 & 1.3 & 3.5 & 0.19 & 0.08 & 0.13 & 0.18 & 0.23 & 0.31\\
Niger & 0.3 & 0.0 & 0.1 & 0.1 & 0.2 & 1.2 & 0.09 & 0.04 & 0.06 & 0.07 & 0.09 & 0.20\\
Nigeria & 1.2 & 0.3 & 0.7 & 1.1 & 1.6 & 2.1 & 0.46 & 0.29 & 0.37 & 0.47 & 0.54 & 0.61\\
Norway & 15.4 & 12.4 & 14.9 & 15.2 & 16.3 & 18.3 & 0.33 & 0.45 & 0.37 & 0.32 & 0.28 & 0.24\\
Pakistan & 2.4 & 0.9 & 1.4 & 2.0 & 2.8 & 4.7 & 0.60 & 0.40 & 0.50 & 0.60 & 0.69 & 0.80\\
Paraguay & 2.5 & 0.9 & 1.9 & 2.5 & 3.0 & 4.5 & 0.35 & 0.33 & 0.40 & 0.36 & 0.33 & 0.32\\
Peru & 1.4 & 0.5 & 1.0 & 1.3 & 1.6 & 2.4 & 0.32 & 0.35 & 0.34 & 0.31 & 0.30 & 0.28\\
Philippines & 2.3 & 0.6 & 1.1 & 1.8 & 2.9 & 4.9 & 0.45 & 0.31 & 0.38 & 0.46 & 0.52 & 0.56\\
Poland & 12.9 & 7.4 & 10.6 & 12.2 & 14.7 & 19.4 & 1.09 & 1.02 & 1.16 & 1.18 & 1.16 & 0.93\\
Portugal & 6.7 & 4.3 & 5.5 & 6.5 & 7.6 & 9.5 & 0.40 & 0.48 & 0.43 & 0.40 & 0.38 & 0.33\\
Romania & 5.2 & 2.1 & 3.6 & 5.2 & 6.8 & 8.2 & 0.61 & 0.62 & 0.70 & 0.67 & 0.58 & 0.49\\
Russia & 8.5 & 4.2 & 6.4 & 7.4 & 9.1 & 15.6 & 1.28 & 1.36 & 1.32 & 1.22 & 1.24 & 1.28\\
Rwanda & 0.1 & 0.0 & 0.0 & 0.0 & 0.1 & 0.6 & 0.04 & 0.02 & 0.02 & 0.02 & 0.03 & 0.10\\
Senegal & 1.9 & 0.5 & 0.9 & 1.7 & 2.3 & 4.0 & 0.23 & 0.13 & 0.16 & 0.24 & 0.27 & 0.35\\
Serbia & 5.4 & 2.5 & 4.4 & 5.4 & 6.5 & 8.1 & 0.64 & 0.51 & 0.63 & 0.66 & 0.68 & 0.74\\
Slovakia & 8.3 & 7.3 & 7.6 & 8.5 & 8.2 & 9.7 & 0.71 & 0.85 & 0.73 & 0.76 & 0.66 & 0.54\\
South Africa & 13.8 & 3.4 & 5.5 & 8.0 & 14.0 & 38.2 & 2.03 & 2.06 & 1.98 & 2.02 & 2.06 & 2.02\\
Spain & 9.9 & 5.4 & 8.2 & 10.3 & 11.9 & 13.9 & 0.42 & 0.44 & 0.44 & 0.43 & 0.42 & 0.39\\
Suriname & 1.0 & 0.4 & 0.7 & 0.9 & 1.3 & 1.6 & 0.13 & 0.14 & 0.14 & 0.14 & 0.14 & 0.12\\
Sweden & 5.7 & 4.7 & 5.7 & 5.4 & 5.8 & 6.9 & 0.19 & 0.22 & 0.23 & 0.18 & 0.18 & 0.17\\
Switzerland & 11.7 & 9.8 & 11.0 & 11.6 & 12.3 & 14.0 & 0.15 & 0.16 & 0.15 & 0.15 & 0.15 & 0.14\\
Taiwan & 24.0 & 15.5 & 20.9 & 24.0 & 27.5 & 31.9 & 1.17 & 1.17 & 1.17 & 1.18 & 1.18 & 1.17\\
Thailand & 6.2 & 1.6 & 3.2 & 5.3 & 8.0 & 12.8 & 1.56 & 1.42 & 1.64 & 1.70 & 1.59 & 1.46\\
Togo & 0.9 & 0.2 & 0.4 & 0.6 & 1.0 & 2.1 & 0.26 & 0.16 & 0.23 & 0.23 & 0.28 & 0.42\\
Turkey & 11.4 & 6.6 & 9.6 & 11.4 & 12.6 & 17.0 & 1.25 & 1.32 & 1.43 & 1.32 & 1.19 & 1.01\\
Uganda & 0.3 & 0.1 & 0.1 & 0.1 & 0.2 & 0.9 & 0.20 & 0.28 & 0.17 & 0.14 & 0.14 & 0.25\\
United Kingdom & 16.3 & 9.5 & 13.8 & 17.2 & 18.5 & 22.8 & 0.55 & 0.68 & 0.61 & 0.56 & 0.49 & 0.40\\
United States & 34.7 & 24.3 & 31.6 & 34.3 & 37.7 & 45.6 & 0.86 & 1.00 & 0.97 & 0.86 & 0.80 & 0.68\\
Uruguay & 2.6 & 1.2 & 1.8 & 2.4 & 3.1 & 4.5 & 0.13 & 0.15 & 0.13 & 0.13 & 0.12 & 0.11\\
Vietnam & 1.1 & 0.4 & 0.7 & 1.0 & 1.4 & 2.0 & 0.51 & 0.56 & 0.53 & 0.51 & 0.49 & 0.44\\*
\end{longtable}
\end{ThreePartTable}
\endgroup{}

\clearpage

% Table C8.2: Sectoral shares and means
\begingroup\fontsize{9}{11}\selectfont

\begin{ThreePartTable}
\begin{TableNotes}
\item \textit{Note: } 
\item This table shows the share of different consumption sectors on the carbon intensity of consumption for all countries of our sample. Values show household-weighted averages and household-weighted standard deviations in parentheses. \textit{Cooking energy} comprises other forms of residential energy excluding electricity.
\end{TableNotes}
\begin{longtable}[t]{l|ccccccl|ccccccl|ccccccl|ccccccl|ccccccl|ccccccl|cccccc}
\caption{\label{tab:A8_Shares}Share of consumption items on carbon intensity of consumption}\\
\toprule
Country & Electricity & Cooking energy & Transport fuels & Food & Goods & Services\\
\midrule
\endfirsthead
\caption[]{Share of consumption items on carbon intensity of consumption \textit{(continued)}}\\
\toprule
Country & Electricity & Cooking energy & Transport fuels & Food & Goods & Services\\
\midrule
\endhead

\endfoot
\bottomrule
\insertTableNotes
\endlastfoot
Argentina & 0.31 & 0.3 & 0.14 & 0.12 & 0.06 & 0.07\\
\nopagebreak
 & (0.21) & (0.2) & (0.21) & (0.12) & (0.07) & (0.09)\\
Armenia & 0.14 & 0.57 & 0.13 & 0.02 & 0 & 0.06\\
\nopagebreak
 & (0.17) & (0.3) & (0.23) & (0.04) & (0.01) & (0.1)\\
Australia & 0.25 & 0.11 & 0.19 & 0.11 & 0.08 & 0.25\\
\nopagebreak
 & (0.17) & (0.14) & (0.16) & (0.06) & (0.08) & (0.18)\\
Austria & 0.06 & 0.24 & 0.36 & 0.05 & 0.07 & 0.21\\
\nopagebreak
 & (0.08) & (0.26) & (0.31) & (0.06) & (0.09) & (0.17)\\
Bangladesh & 0.11 & 0.21 & 0.01 & 0.16 & 0.27 & 0.23\\
\nopagebreak
 & (0.13) & (0.16) & (0.08) & (0.09) & (0.12) & (0.16)\\
Barbados & 0.17 & 0.18 & 0.31 & 0.08 & 0.09 & 0.17\\
\nopagebreak
 & (0.17) & (0.23) & (0.36) & (0.11) & (0.11) & (0.18)\\
Belgium & 0.06 & 0.38 & 0.24 & 0.04 & 0.05 & 0.23\\
\nopagebreak
 & (0.08) & (0.24) & (0.2) & (0.03) & (0.05) & (0.17)\\
Benin & 0.01 & 0.18 & 0.43 & 0.03 & 0.02 & 0.33\\
\nopagebreak
 & (0.03) & (0.3) & (0.45) & (0.07) & (0.05) & (0.36)\\
Bolivia & 0.32 & 0.13 & 0.1 & 0.12 & 0.02 & 0.31\\
\nopagebreak
 & (0.18) & (0.11) & (0.19) & (0.09) & (0.03) & (0.2)\\
Brazil & 0.12 & 0.2 & 0.19 & 0.15 & 0.15 & 0.2\\
\nopagebreak
 & (0.12) & (0.19) & (0.26) & (0.14) & (0.15) & (0.18)\\
Bulgaria & 0.48 & 0.05 & 0.24 & 0.08 & 0.03 & 0.11\\
\nopagebreak
 & (0.24) & (0.17) & (0.25) & (0.06) & (0.03) & (0.11)\\
Burkina Faso & 0.01 & 0.03 & 0.53 & 0.08 & 0.16 & 0.18\\
\nopagebreak
 & (0.07) & (0.11) & (0.46) & (0.13) & (0.23) & (0.25)\\
Cambodia & 0.13 & 0.11 & 0.5 & 0.08 & 0.04 & 0.14\\
\nopagebreak
 & (0.14) & (0.14) & (0.3) & (0.09) & (0.07) & (0.14)\\
Canada & 0.11 & 0.19 & 0.23 & 0.09 & 0.11 & 0.27\\
\nopagebreak
 & (0.11) & (0.24) & (0.19) & (0.08) & (0.11) & (0.19)\\
Chile & 0.17 & 0.11 & 0.13 & 0.19 & 0.12 & 0.28\\
\nopagebreak
 & (0.12) & (0.13) & (0.16) & (0.1) & (0.09) & (0.15)\\
Colombia & 0.09 & 0.31 & 0.11 & 0.11 & 0.18 & 0.2\\
\nopagebreak
 & (0.12) & (0.29) & (0.2) & (0.14) & (0.2) & (0.19)\\
Costa Rica & 0.01 & 0.14 & 0.32 & 0.19 & 0.03 & 0.31\\
\nopagebreak
 & (0.02) & (0.23) & (0.39) & (0.22) & (0.05) & (0.3)\\
Croatia & 0.17 & 0.26 & 0.37 & 0.04 & 0.01 & 0.15\\
\nopagebreak
 & (0.18) & (0.29) & (0.31) & (0.06) & (0.02) & (0.18)\\
Cyprus & 0.15 & 0.14 & 0.53 & 0.03 & 0 & 0.15\\
\nopagebreak
 & (0.13) & (0.16) & (0.28) & (0.04) & (0.01) & (0.16)\\
Czechia & 0.22 & 0.23 & 0.23 & 0.03 & 0.13 & 0.17\\
\nopagebreak
 & (0.14) & (0.25) & (0.19) & (0.01) & (0.11) & (0.13)\\
Côte d’Ivoire & 0.17 & 0.16 & 0.2 & 0.11 & 0.1 & 0.25\\
\nopagebreak
 & (0.23) & (0.26) & (0.36) & (0.15) & (0.13) & (0.23)\\
Denmark & 0.08 & 0.11 & 0.35 & 0.06 & 0.09 & 0.31\\
\nopagebreak
 & (0.06) & (0.21) & (0.25) & (0.04) & (0.08) & (0.18)\\
Dominican Republic & 0.1 & 0.33 & 0.23 & 0.11 & 0.04 & 0.19\\
\nopagebreak
 & (0.11) & (0.23) & (0.31) & (0.11) & (0.07) & (0.18)\\
Ecuador & 0.06 & 0.41 & 0.19 & 0.18 & 0.07 & 0.09\\
\nopagebreak
 & (0.08) & (0.28) & (0.32) & (0.15) & (0.08) & (0.12)\\
Egypt & 0.18 & 0.31 & 0.08 & 0.08 & 0.11 & 0.24\\
\nopagebreak
 & (0.06) & (0.1) & (0.19) & (0.04) & (0.04) & (0.12)\\
El Salvador & 0.22 & 0.5 & 0.1 & 0.02 & 0.06 & 0.1\\
\nopagebreak
 & (0.25) & (0.38) & (0.24) & (0.07) & (0.13) & (0.16)\\
Estonia & 0.37 & 0.03 & 0.37 & 0.05 & 0.03 & 0.14\\
\nopagebreak
 & (0.27) & (0.1) & (0.3) & (0.09) & (0.04) & (0.15)\\
Ethiopia & 0 & 0.14 & 0.14 & 0.49 & 0.18 & 0.05\\
\nopagebreak
 & (0) & (0.24) & (0.21) & (0.25) & (0.15) & (0.07)\\
Finland & 0.09 & 0.02 & 0.2 & 0.03 & 0.13 & 0.52\\
\nopagebreak
 & (0.08) & (0.1) & (0.22) & (0.02) & (0.1) & (0.2)\\
France & 0.11 & 0.19 & 0.23 & 0.09 & 0.09 & 0.3\\
\nopagebreak
 & (0.12) & (0.26) & (0.27) & (0.09) & (0.1) & (0.23)\\
Georgia & 0.09 & 0.64 & 0.18 & 0.02 & 0 & 0.06\\
\nopagebreak
 & (0.14) & (0.31) & (0.3) & (0.07) & (0.03) & (0.11)\\
Germany & 0.1 & 0.33 & 0.23 & 0.07 & 0.12 & 0.14\\
\nopagebreak
 & (0.08) & (0.28) & (0.21) & (0.05) & (0.1) & (0.1)\\
Ghana & 0.3 & 0.11 & 0.07 & 0.22 & 0.08 & 0.23\\
\nopagebreak
 & (0.27) & (0.22) & (0.23) & (0.2) & (0.11) & (0.22)\\
Greece & 0.28 & 0.15 & 0.22 & 0.06 & 0.03 & 0.26\\
\nopagebreak
 & (0.17) & (0.17) & (0.2) & (0.03) & (0.02) & (0.15)\\
Guatemala & 0.14 & 0.3 & 0.16 & 0.07 & 0.06 & 0.26\\
\nopagebreak
 & (0.19) & (0.3) & (0.3) & (0.1) & (0.09) & (0.29)\\
Guinea-Bissau & 0.03 & 0.02 & 0.1 & 0.19 & 0.17 & 0.49\\
\nopagebreak
 & (0.09) & (0.1) & (0.27) & (0.18) & (0.17) & (0.31)\\
Hungary & 0.13 & 0.44 & 0.23 & 0.03 & 0.03 & 0.14\\
\nopagebreak
 & (0.12) & (0.28) & (0.24) & (0.03) & (0.04) & (0.13)\\
India & 0.13 & 0.13 & 0.16 & 0.21 & 0.26 & 0.11\\
\nopagebreak
 & (0.08) & (0.1) & (0.17) & (0.11) & (0.12) & (0.06)\\
Indonesia & 0.29 & 0.15 & 0.39 & 0.06 & 0.06 & 0.05\\
\nopagebreak
 & (0.18) & (0.13) & (0.25) & (0.05) & (0.06) & (0.06)\\
Iraq & 0.37 & 0.26 & 0.15 & 0.05 & 0.05 & 0.12\\
\nopagebreak
 & (0.23) & (0.18) & (0.24) & (0.05) & (0.07) & (0.11)\\
Ireland & 0.15 & 0.29 & 0.25 & 0.05 & 0.08 & 0.17\\
\nopagebreak
 & (0.12) & (0.23) & (0.21) & (0.04) & (0.06) & (0.14)\\
Israel & 0.42 & 0.08 & 0.14 & 0.05 & 0.06 & 0.26\\
\nopagebreak
 & (0.2) & (0.09) & (0.19) & (0.05) & (0.07) & (0.16)\\
Italy & 0.11 & 0.25 & 0.26 & 0.11 & 0.1 & 0.17\\
\nopagebreak
 & (0.09) & (0.19) & (0.2) & (0.07) & (0.09) & (0.13)\\
Jordan & 0.1 & 0.34 & 0.4 & 0.05 & 0.01 & 0.1\\
\nopagebreak
 & (0.07) & (0.25) & (0.36) & (0.04) & (0.01) & (0.11)\\
Kenya & 0.02 & 0.34 & 0.03 & 0.09 & 0.18 & 0.33\\
\nopagebreak
 & (0.06) & (0.28) & (0.15) & (0.13) & (0.19) & (0.3)\\
Latvia & 0.04 & 0.48 & 0.28 & 0.01 & 0.01 & 0.18\\
\nopagebreak
 & (0.05) & (0.34) & (0.29) & (0.02) & (0.03) & (0.2)\\
Liberia & 0.02 & 0 & 0.04 & 0.35 & 0.12 & 0.47\\
\nopagebreak
 & (0.06) & (0.04) & (0.18) & (0.28) & (0.11) & (0.33)\\
Lithuania & 0.03 & 0.33 & 0.34 & 0.07 & 0.09 & 0.13\\
\nopagebreak
 & (0.06) & (0.35) & (0.38) & (0.12) & (0.13) & (0.19)\\
Luxembourg & 0 & 0.65 & 0.25 & 0 & 0.02 & 0.07\\
\nopagebreak
 & (0.01) & (0.27) & (0.26) & (0) & (0.03) & (0.12)\\
Malawi & 0.02 & 0.01 & 0.02 & 0.42 & 0.2 & 0.32\\
\nopagebreak
 & (0.1) & (0.06) & (0.14) & (0.27) & (0.17) & (0.33)\\
Maldives & 0.25 & 0.28 & 0.13 & 0.02 & 0.05 & 0.28\\
\nopagebreak
 & (0.17) & (0.21) & (0.15) & (0.04) & (0.05) & (0.23)\\
Mali & 0.01 & 0.02 & 0.52 & 0.1 & 0.09 & 0.26\\
\nopagebreak
 & (0.05) & (0.1) & (0.47) & (0.17) & (0.15) & (0.32)\\
Mexico & 0.12 & 0.18 & 0.23 & 0.15 & 0.15 & 0.17\\
\nopagebreak
 & (0.13) & (0.2) & (0.29) & (0.12) & (0.12) & (0.17)\\
Mongolia & 0.02 & 0.41 & 0.34 & 0.02 & 0.01 & 0.2\\
\nopagebreak
 & (0.05) & (0.38) & (0.35) & (0.04) & (0.03) & (0.27)\\
Morocco & 0.19 & 0.3 & 0.1 & 0.18 & 0.05 & 0.18\\
\nopagebreak
 & (0.11) & (0.15) & (0.2) & (0.09) & (0.05) & (0.11)\\
Mozambique & 0 & 0.03 & 0.12 & 0.23 & 0.11 & 0.5\\
\nopagebreak
 & (0.03) & (0.11) & (0.3) & (0.3) & (0.22) & (0.38)\\
Myanmar (Burma) & 0.08 & 0.06 & 0.11 & 0.19 & 0.12 & 0.45\\
\nopagebreak
 & (0.14) & (0.13) & (0.24) & (0.17) & (0.19) & (0.29)\\
Netherlands & 0.15 & 0.27 & 0.17 & 0.04 & 0.06 & 0.31\\
\nopagebreak
 & (0.08) & (0.12) & (0.14) & (0.03) & (0.05) & (0.14)\\
Nicaragua & 0.09 & 0.27 & 0.1 & 0.09 & 0.04 & 0.41\\
\nopagebreak
 & (0.16) & (0.34) & (0.26) & (0.17) & (0.09) & (0.35)\\
Niger & 0.06 & 0 & 0.11 & 0.02 & 0.52 & 0.29\\
\nopagebreak
 & (0.17) & (0.02) & (0.29) & (0.03) & (0.33) & (0.28)\\
Nigeria & 0.11 & 0.17 & 0.18 & 0.07 & 0.09 & 0.38\\
\nopagebreak
 & (0.18) & (0.21) & (0.28) & (0.12) & (0.13) & (0.28)\\
Norway & 0.02 & 0.31 & 0.27 & 0.04 & 0.09 & 0.27\\
\nopagebreak
 & (0.04) & (0.31) & (0.26) & (0.05) & (0.1) & (0.24)\\
Pakistan & 0.3 & 0.2 & 0.19 & 0.1 & 0.16 & 0.06\\
\nopagebreak
 & (0.21) & (0.22) & (0.25) & (0.1) & (0.13) & (0.06)\\
Paraguay & 0 & 0.34 & 0.35 & 0.11 & 0.01 & 0.18\\
\nopagebreak
 & (0.03) & (0.31) & (0.35) & (0.21) & (0.04) & (0.24)\\
Peru & 0.25 & 0.34 & 0.09 & 0.14 & 0.06 & 0.12\\
\nopagebreak
 & (0.18) & (0.23) & (0.19) & (0.12) & (0.08) & (0.16)\\
Philippines & 0.39 & 0.14 & 0.11 & 0.15 & 0.03 & 0.17\\
\nopagebreak
 & (0.2) & (0.14) & (0.18) & (0.11) & (0.04) & (0.12)\\
Poland & 0.2 & 0.18 & 0.23 & 0.08 & 0.05 & 0.26\\
\nopagebreak
 & (0.18) & (0.25) & (0.19) & (0.06) & (0.05) & (0.17)\\
Portugal & 0.21 & 0.17 & 0.32 & 0.04 & 0.04 & 0.21\\
\nopagebreak
 & (0.13) & (0.13) & (0.22) & (0.03) & (0.03) & (0.15)\\
Romania & 0.22 & 0.41 & 0.17 & 0.04 & 0.03 & 0.13\\
\nopagebreak
 & (0.19) & (0.27) & (0.24) & (0.05) & (0.05) & (0.12)\\
Russia & 0 & 0.03 & 0.18 & 0.37 & 0.2 & 0.22\\
\nopagebreak
 & (0) & (0.14) & (0.26) & (0.23) & (0.16) & (0.2)\\
Rwanda & 0.09 & 0.06 & 0.02 & 0.35 & 0.13 & 0.35\\
\nopagebreak
 & (0.2) & (0.2) & (0.13) & (0.26) & (0.12) & (0.3)\\
Senegal & 0.2 & 0.2 & 0.08 & 0.06 & 0.07 & 0.39\\
\nopagebreak
 & (0.22) & (0.24) & (0.23) & (0.08) & (0.09) & (0.27)\\
Serbia & 0.39 & 0.17 & 0.12 & 0.09 & 0.1 & 0.13\\
\nopagebreak
 & (0.21) & (0.26) & (0.16) & (0.06) & (0.08) & (0.12)\\
Slovakia & 0.14 & 0.37 & 0.3 & 0.01 & 0.04 & 0.13\\
\nopagebreak
 & (0.14) & (0.33) & (0.25) & (0.01) & (0.1) & (0.15)\\
South Africa & 0.3 & 0 & 0.13 & 0.16 & 0.15 & 0.25\\
\nopagebreak
 & (0.21) & (0.04) & (0.21) & (0.14) & (0.16) & (0.17)\\
Spain & 0.18 & 0.13 & 0.24 & 0.1 & 0.11 & 0.25\\
\nopagebreak
 & (0.14) & (0.16) & (0.24) & (0.08) & (0.12) & (0.16)\\
Suriname & 0.11 & 0.19 & 0.09 & 0.11 & 0.05 & 0.45\\
\nopagebreak
 & (0.14) & (0.24) & (0.21) & (0.15) & (0.11) & (0.28)\\
Sweden & 0.05 & 0.04 & 0.48 & 0.04 & 0.13 & 0.27\\
\nopagebreak
 & (0.07) & (0.15) & (0.34) & (0.07) & (0.16) & (0.26)\\
Switzerland & 0.01 & 0.06 & 0.3 & 0.04 & 0.14 & 0.45\\
\nopagebreak
 & (0.02) & (0.17) & (0.25) & (0.05) & (0.17) & (0.23)\\
Taiwan & 0.1 & 0.06 & 0.32 & 0.05 & 0.11 & 0.35\\
\nopagebreak
 & (0.05) & (0.04) & (0.2) & (0.03) & (0.08) & (0.17)\\
Thailand & 0.17 & 0.05 & 0.42 & 0.07 & 0.18 & 0.11\\
\nopagebreak
 & (0.15) & (0.11) & (0.26) & (0.09) & (0.17) & (0.15)\\
Togo & 0.06 & 0.08 & 0.27 & 0.07 & 0.03 & 0.5\\
\nopagebreak
 & (0.14) & (0.25) & (0.43) & (0.12) & (0.06) & (0.4)\\
Turkey & 0.16 & 0.36 & 0.09 & 0.07 & 0.15 & 0.17\\
\nopagebreak
 & (0.15) & (0.29) & (0.14) & (0.07) & (0.14) & (0.15)\\
Uganda & 0.02 & 0.3 & 0.05 & 0.2 & 0.14 & 0.28\\
\nopagebreak
 & (0.08) & (0.37) & (0.21) & (0.23) & (0.18) & (0.31)\\
United Kingdom & 0.08 & 0.41 & 0.17 & 0.07 & 0.07 & 0.19\\
\nopagebreak
 & (0.11) & (0.3) & (0.2) & (0.08) & (0.1) & (0.19)\\
United States & 0.2 & 0.2 & 0.32 & 0.07 & 0.07 & 0.14\\
\nopagebreak
 & (0.15) & (0.21) & (0.19) & (0.08) & (0.11) & (0.13)\\
Uruguay & 0.05 & 0.21 & 0.37 & 0.1 & 0.03 & 0.24\\
\nopagebreak
 & (0.06) & (0.21) & (0.33) & (0.09) & (0.05) & (0.21)\\
Vietnam & 0.24 & 0.05 & 0.07 & 0.22 & 0.18 & 0.24\\
\nopagebreak
 & (0.12) & (0.06) & (0.06) & (0.1) & (0.11) & (0.13)\\*
\end{longtable}
\end{ThreePartTable}
\endgroup{}

\clearpage

% Table C8: Theil index and robustness checks
\begingroup\fontsize{11}{13}\selectfont

\begin{ThreePartTable}
\begin{TableNotes}
\item \textit{Note: } 
\item This table shows the share of heterogeneity in carbon intensity of consumption that can be explained with heterogeneity within expenditure groups in comparison to heterogeneity between expenditure groups in \%. Column (I) uses the Theil index ($\alpha = 1$) with expenditure quintiles. Column (II) uses the Theil index ($\alpha = 2$) with expenditure quintiles which accounts for greater overall heterogeneity and longer right tails. Column (III) uses the Theil index ($\alpha = 1$) with expenditure deciles. Column (IV) uses the Theil index ($\alpha = 1$) with expenditure percentiles.
\end{TableNotes}
\begin{longtable}[t]{l|ccccl|ccccl|ccccl|ccccl|cccc}
\caption{\label{tab:ATBD2}Within-group heterogeneity as a share of heterogeneity in carbon intensity of consumption}\\
\toprule
Country & (I) & (II) & (III) & (IV)\\
\midrule
\endfirsthead
\caption[]{Within-group heterogeneity as a share of heterogeneity in carbon intensity of consumption \textit{(continued)}}\\
\toprule
Country & (I) & (II) & (III) & (IV)\\
\midrule
\endhead

\endfoot
\bottomrule
\insertTableNotes
\endlastfoot
Argentina & 84\% & 87\% & 83\% & 82\%\\
Armenia & 96\% & 97\% & 95\% & 93\%\\
Australia & 83\% & 85\% & 81\% & 80\%\\
Austria & 93\% & 94\% & 93\% & 91\%\\
Bangladesh & 98\% & 98\% & 97\% & 95\%\\
Barbados & 99\% & 99\% & 99\% & 93\%\\
Belgium & 96\% & 96\% & 95\% & 93\%\\
Benin & 98\% & 98\% & 97\% & 95\%\\
Bolivia & 99\% & 99\% & 98\% & 97\%\\
Brazil & 93\% & 95\% & 92\% & 90\%\\
Bulgaria & 100\% & 100\% & 99\% & 93\%\\
Burkina Faso & 98\% & 98\% & 98\% & 95\%\\
Cambodia & 99\% & 99\% & 99\% & 89\%\\
Canada & 97\% & 98\% & 97\% & 95\%\\
Chile & 78\% & 81\% & 76\% & 75\%\\
Colombia & 93\% & 95\% & 92\% & 92\%\\
Costa Rica & 99\% & 100\% & 99\% & 98\%\\
Croatia & 99\% & 99\% & 97\% & 90\%\\
Cyprus & 96\% & 97\% & 96\% & 92\%\\
Czechia & 98\% & 98\% & 98\% & 95\%\\
Côte d’Ivoire & 99\% & 99\% & 99\% & 98\%\\
Denmark & 98\% & 99\% & 98\% & 94\%\\
Dominican Republic & 97\% & 97\% & 97\% & 95\%\\
Ecuador & 99\% & 99\% & 98\% & 96\%\\
Egypt & 98\% & 98\% & 97\% & 96\%\\
El Salvador & 92\% & 93\% & 92\% & 90\%\\
Estonia & 96\% & 96\% & 95\% & 92\%\\
Ethiopia & 96\% & 98\% & 96\% & 93\%\\
Finland & 97\% & 97\% & 96\% & 94\%\\
France & 99\% & 99\% & 99\% & 98\%\\
Georgia & 100\% & 100\% & 99\% & 98\%\\
Germany & 95\% & 96\% & 95\% & 95\%\\
Ghana & 95\% & 97\% & 94\% & 92\%\\
Greece & 94\% & 95\% & 93\% & 92\%\\
Guatemala & 79\% & 84\% & 78\% & 77\%\\
Guinea-Bissau & 95\% & 97\% & 94\% & 91\%\\
Hungary & 99\% & 99\% & 99\% & 97\%\\
India & 95\% & 95\% & 95\% & 94\%\\
Indonesia & 99\% & 99\% & 98\% & 98\%\\
Iraq & 91\% & 93\% & 90\% & 86\%\\
Ireland & 94\% & 95\% & 93\% & 91\%\\
Israel & 80\% & 83\% & 78\% & 75\%\\
Italy & 94\% & 95\% & 94\% & 93\%\\
Jordan & 97\% & 98\% & 97\% & 95\%\\
Kenya & 98\% & 98\% & 98\% & 97\%\\
Latvia & 92\% & 93\% & 91\% & 88\%\\
Liberia & 90\% & 94\% & 89\% & 87\%\\
Lithuania & 99\% & 99\% & 99\% & 95\%\\
Luxembourg & 89\% & 91\% & 88\% & 85\%\\
Malawi & 86\% & 97\% & 84\% & 80\%\\
Maldives & 98\% & 98\% & 98\% & 95\%\\
Mali & 99\% & 99\% & 99\% & 97\%\\
Mexico & 99\% & 99\% & 99\% & 99\%\\
Mongolia & 97\% & 98\% & 96\% & 94\%\\
Morocco & 98\% & 98\% & 97\% & 96\%\\
Mozambique & 93\% & 96\% & 90\% & 85\%\\
Myanmar (Burma) & 93\% & 95\% & 93\% & 89\%\\
Netherlands & 89\% & 90\% & 89\% & 88\%\\
Nicaragua & 88\% & 91\% & 87\% & 85\%\\
Niger & 80\% & 89\% & 78\% & 74\%\\
Nigeria & 90\% & 91\% & 89\% & 89\%\\
Norway & 92\% & 94\% & 91\% & 89\%\\
Pakistan & 86\% & 87\% & 85\% & 85\%\\
Paraguay & 99\% & 100\% & 99\% & 98\%\\
Peru & 99\% & 99\% & 99\% & 99\%\\
Philippines & 81\% & 83\% & 80\% & 80\%\\
Poland & 99\% & 99\% & 98\% & 98\%\\
Portugal & 89\% & 90\% & 87\% & 86\%\\
Romania & 96\% & 98\% & 96\% & 95\%\\
Russia & 99\% & 100\% & 99\% & 97\%\\
Rwanda & 85\% & 95\% & 81\% & 77\%\\
Senegal & 88\% & 92\% & 88\% & 86\%\\
Serbia & 98\% & 99\% & 98\% & 94\%\\
Slovakia & 97\% & 98\% & 96\% & 94\%\\
South Africa & 100\% & 100\% & 100\% & 99\%\\
Spain & 99\% & 99\% & 99\% & 98\%\\
Suriname & 99\% & 99\% & 99\% & 92\%\\
Sweden & 99\% & 99\% & 99\% & 94\%\\
Switzerland & 99\% & 99\% & 99\% & 98\%\\
Taiwan & 100\% & 100\% & 100\% & 99\%\\
Thailand & 98\% & 98\% & 97\% & 95\%\\
Togo & 96\% & 97\% & 94\% & 91\%\\
Turkey & 97\% & 98\% & 96\% & 95\%\\
Uganda & 96\% & 98\% & 95\% & 93\%\\
United Kingdom & 95\% & 96\% & 94\% & 92\%\\
United States & 92\% & 93\% & 92\% & 89\%\\
Uruguay & 98\% & 99\% & 98\% & 96\%\\
Vietnam & 93\% & 94\% & 93\% & 92\%\\*
\end{longtable}
\end{ThreePartTable}
\endgroup{}

\clearpage

% Feature overview
\input{2_Tables/Table_Feature_Overview}
\clearpage

% Table C8: Hyperparameters
\begingroup\fontsize{8}{10}\selectfont

\begin{ThreePartTable}
\begin{TableNotes}
\item \textit{Note: } 
\item This table shows hyperparameters selected for fitting boosted regression tree models after hyperparameter tuning. \texttt{max\_depth} is the maximum depth of trees; $\eta$ is the learning rate; \texttt{mtry} is the number of features included in each tree. MAE is the mean absolute error of predictions; RMSE is the root mean squared error of predictions; R\textsuperscript{2} is the squared correlation of prediction errors and a measure for goodness of fit. Unit of MAE and RMSE is kgCO\textsubscript{2} per USD. We show MAE, RMSE and R\textsuperscript{2} for five-fold cross-validation on the entire dataset. Note that we repeat five-fold cross-validation with selected hyperparameters, e.g. in Table \ref{tab:A8}. We include survey weights in all analyses.
\end{TableNotes}
\begin{longtable}[t]{l|rcr|cccl|rcr|cccl|rcr|cccl|rcr|cccl|rcr|cccl|rcr|cccl|rcr|ccc}
\caption{\label{tab:C9}Hyperparameters for boosted regression tree models}\\
\toprule
\multicolumn{1}{c}{ } & \multicolumn{3}{c}{Selected hyperparameters} & \multicolumn{3}{c}{Model performance} \\
\cmidrule(l{3pt}r{3pt}){2-4} \cmidrule(l{3pt}r{3pt}){5-7}
Country & \texttt{max\_depth} & $\eta$ & \texttt{mtry} & MAE & RMSE & R\textsuperscript{2}\\
\midrule
\endfirsthead
\caption[]{Hyperparameters for boosted regression tree models \textit{(continued)}}\\
\toprule
Country & \texttt{max\_depth} & $\eta$ & \texttt{mtry} & MAE & RMSE & R\textsuperscript{2}\\
\midrule
\endhead

\endfoot
\bottomrule
\insertTableNotes
\endlastfoot
Argentina & 4 & 0.0202 & 31 & 0.50 & 0.75 & 0.29\\
Armenia & 5 & 0.0069 & 18 & 0.43 & 0.63 & 0.26\\
Australia & 3 & 0.0110 & 13 & 0.18 & 0.25 & 0.24\\
Austria & 3 & 0.0110 & 16 & 0.12 & 0.16 & 0.20\\
Bangladesh & 5 & 0.0095 & 10 & 0.08 & 0.13 & 0.17\\
Barbados & 3 & 0.0063 & 25 & 0.26 & 0.40 & 0.23\\
Belgium & 3 & 0.0146 & 10 & 0.14 & 0.20 & 0.11\\
Benin & 5 & 0.0095 & 23 & 0.22 & 0.35 & 0.34\\
Bolivia & 5 & 0.0095 & 16 & 0.11 & 0.15 & 0.38\\
Brazil & 6 & 0.0031 & 24 & 0.18 & 0.31 & 0.15\\
Bulgaria & 3 & 0.0028 & 5 & 0.22 & 0.40 & 0.01\\
Burkina Faso & 5 & 0.0069 & 26 & 0.19 & 0.31 & 0.46\\
Cambodia & 4 & 0.0047 & 22 & 0.21 & 0.31 & 0.20\\
Canada & 5 & 0.0095 & 8 & 0.18 & 0.28 & 0.14\\
Chile & 3 & 0.0063 & 6 & 0.12 & 0.18 & 0.23\\
Colombia & 4 & 0.0202 & 13 & 0.23 & 0.37 & 0.17\\
Costa Rica & 3 & 0.0110 & 15 & 0.12 & 0.18 & 0.23\\
Croatia & 3 & 0.0028 & 5 & 0.22 & 0.33 & 0.07\\
Cyprus & 3 & 0.0028 & 8 & 0.10 & 0.14 & 0.08\\
Czechia & 3 & 0.0028 & 4 & 0.32 & 0.48 & 0.09\\
Côte d’Ivoire & 9 & 0.0066 & 32 & 0.13 & 0.21 & 0.41\\
Denmark & 3 & 0.0028 & 6 & 0.08 & 0.11 & 0.05\\
Dominican Republic & 5 & 0.0095 & 25 & 0.11 & 0.17 & 0.42\\
Ecuador & 5 & 0.0087 & 20 & 0.14 & 0.31 & 0.44\\
Egypt & 3 & 0.0110 & 13 & 0.21 & 0.32 & 0.38\\
El Salvador & 5 & 0.0087 & 24 & 0.18 & 0.30 & 0.33\\
Estonia & 4 & 0.0019 & 8 & 0.22 & 0.29 & 0.06\\
Ethiopia & 5 & 0.0095 & 17 & 0.03 & 0.07 & 0.15\\
Finland & 3 & 0.0028 & 7 & 0.12 & 0.15 & 0.10\\
France & 5 & 0.0027 & 30 & 0.13 & 0.22 & 0.23\\
Georgia & 9 & 0.0066 & 14 & 0.49 & 0.72 & 0.29\\
Germany & 5 & 0.0095 & 16 & 0.14 & 0.20 & 0.34\\
Ghana & 5 & 0.0095 & 21 & 0.13 & 0.26 & 0.35\\
Greece & 4 & 0.0047 & 6 & 0.11 & 0.14 & 0.10\\
Guatemala & 5 & 0.0095 & 17 & 0.14 & 0.22 & 0.43\\
Guinea-Bissau & 5 & 0.0069 & 26 & 0.14 & 0.24 & 0.17\\
Hungary & 4 & 0.0079 & 11 & 0.30 & 0.39 & 0.03\\
India & 10 & 0.0126 & 36 & 0.20 & 0.27 & 0.48\\
Indonesia & 5 & 0.0569 & 28 & 0.42 & 0.57 & 0.37\\
Iraq & 4 & 0.0079 & 20 & 0.36 & 0.53 & 0.26\\
Ireland & 3 & 0.0045 & 10 & 0.17 & 0.26 & 0.13\\
Israel & 5 & 0.0095 & 20 & 0.12 & 0.17 & 0.38\\
Italy & 4 & 0.0047 & 6 & 0.12 & 0.16 & 0.10\\
Jordan & 3 & 0.0110 & 47 & 0.29 & 0.40 & 0.58\\
Kenya & 9 & 0.0066 & 33 & 0.21 & 0.34 & 0.15\\
Latvia & 3 & 0.0110 & 5 & 0.33 & 0.45 & 0.22\\
Liberia & 5 & 0.0095 & 20 & 0.14 & 0.24 & 0.12\\
Lithuania & 3 & 0.0045 & 9 & 0.29 & 0.44 & 0.06\\
Luxembourg & 4 & 0.0047 & 6 & 0.13 & 0.18 & 0.17\\
Malawi & 3 & 0.0028 & 16 & 0.03 & 0.12 & 0.19\\
Maldives & 9 & 0.0018 & 17 & 0.10 & 0.13 & 0.14\\
Mali & 9 & 0.0066 & 15 & 0.15 & 0.23 & 0.38\\
Mexico & 9 & 0.0066 & 26 & 0.33 & 0.48 & 0.29\\
Mongolia & 6 & 0.0031 & 4 & 0.63 & 0.98 & 0.22\\
Morocco & 3 & 0.0110 & 13 & 0.17 & 0.24 & 0.07\\
Mozambique & 5 & 0.0095 & 16 & 0.22 & 0.41 & 0.17\\
Myanmar (Burma) & 6 & 0.0031 & 21 & 0.25 & 0.37 & 0.15\\
Netherlands & 3 & 0.0063 & 6 & 0.07 & 0.09 & 0.16\\
Nicaragua & 6 & 0.0037 & 21 & 0.09 & 0.16 & 0.50\\
Niger & 3 & 0.0045 & 24 & 0.07 & 0.15 & 0.52\\
Nigeria & 3 & 0.0450 & 37 & 0.19 & 0.28 & 0.35\\
Norway & 3 & 0.0063 & 15 & 0.17 & 0.27 & 0.25\\
Pakistan & 5 & 0.0095 & 7 & 0.24 & 0.32 & 0.24\\
Paraguay & 7 & 0.0036 & 17 & 0.16 & 0.25 & 0.22\\
Peru & 5 & 0.0095 & 17 & 0.11 & 0.18 & 0.31\\
Philippines & 4 & 0.0202 & 17 & 0.12 & 0.17 & 0.45\\
Poland & 6 & 0.0010 & 6 & 0.66 & 1.39 & 0.03\\
Portugal & 3 & 0.0110 & 7 & 0.10 & 0.14 & 0.15\\
Romania & 5 & 0.0095 & 15 & 0.27 & 0.59 & 0.21\\
Russia & 7 & 0.0036 & 32 & 0.41 & 0.69 & 0.27\\
Rwanda & 5 & 0.0095 & 10 & 0.03 & 0.11 & 0.47\\
Senegal & 4 & 0.0047 & 22 & 0.11 & 0.19 & 0.33\\
Serbia & 7 & 0.0012 & 11 & 0.39 & 1.11 & 0.04\\
Slovakia & 4 & 0.0047 & 4 & 0.44 & 0.75 & 0.13\\
South Africa & 5 & 0.0095 & 16 & 0.48 & 0.68 & 0.36\\
Spain & 5 & 0.0069 & 26 & 0.13 & 0.18 & 0.09\\
Suriname & 6 & 0.0010 & 31 & 0.07 & 0.12 & 0.02\\
Sweden & 4 & 0.0015 & 3 & 0.12 & 0.21 & 0.05\\
Switzerland & 6 & 0.0031 & 12 & 0.07 & 0.11 & 0.17\\
Taiwan & 5 & 0.0095 & 8 & 0.20 & 0.25 & 0.36\\
Thailand & 5 & 0.0095 & 13 & 0.42 & 0.56 & 0.33\\
Togo & 5 & 0.0069 & 23 & 0.18 & 0.36 & 0.43\\
Turkey & 7 & 0.0036 & 13 & 0.56 & 0.83 & 0.19\\
Uganda & 5 & 0.0069 & 20 & 0.14 & 0.33 & 0.27\\
United Kingdom & 4 & 0.0047 & 13 & 0.25 & 0.40 & 0.18\\
United States & 3 & 0.0110 & 13 & 0.29 & 0.40 & 0.12\\
Uruguay & 4 & 0.0047 & 13 & 0.05 & 0.08 & 0.33\\
Vietnam & 3 & 0.0450 & 49 & 0.09 & 0.13 & 0.29\\*
\end{longtable}
\end{ThreePartTable}
\endgroup{}

\clearpage

% Table C9: Vertical and horizontal indicators
\begingroup\fontsize{9}{11}\selectfont

\begin{ThreePartTable}
\begin{TableNotes}
\item \textit{Note: } 
\item This table shows the median carbon intensity in the first expenditure quintile ($\overline{e}_{r}^{1}$) and in the fifth quintile ($\overline{e}_{r}^{5}$). It displays the difference between the 5\textsuperscript{th} (20\textsuperscript{th}) and 95\textsuperscript{th} (80\textsuperscript{th}) within-quintile percentile for the first ($\overline{H}_{r}^{1}$ and $\overline{H}_{r}^{1*}$) and the fifth quintile ($\overline{H}_{r}^{5}$ and $\overline{H}_{r}^{5*}$). It also compares the median carbon intensity in the first income quintile to that in the fifth quintile ($\widehat{V}$$_{r}^{1}$). Lastly, it displays our comparison index that enables the comparison of within-quintile variation between the first and fifth quintile ($\widehat{H}_{r}^{1}$ and $\widehat{H}_{r}^{1*}$, respectively).
\end{TableNotes}
\begin{longtable}[t]{l|cc|cccc|cccl|cc|cccc|cccl|cc|cccc|cccl|cc|cccc|cccl|cc|cccc|cccl|cc|cccc|cccl|cc|cccc|cccl|cc|cccc|cccl|cc|cccc|cccl|cc|cccc|ccc}
\caption{\label{tab:A7}Comparing median carbon intensity and horizontal heterogeneity between first and fifth expenditure quintile}\\
\toprule
\multicolumn{1}{c}{Country} & \multicolumn{1}{c}{$\overline{e}_{r}^{1}$} & \multicolumn{1}{c}{$\overline{e}_{r}^{5}$} & \multicolumn{1}{c}{$\overline{H}_{r}^{1}$} & \multicolumn{1}{c}{$\overline{H}_{r}^{5}$} & \multicolumn{1}{c}{$\overline{H}_{r}^{1*}$} & \multicolumn{1}{c}{$\overline{H}_{r}^{5*}$} & \multicolumn{1}{c}{$\widehat{V}_{r}^{1}$} & \multicolumn{1}{c}{$\widehat{H}_{r}^{1}$} & \multicolumn{1}{c}{$\widehat{H}_{r}^{1*}$} \\
\cmidrule(l{3pt}r{3pt}){1-1} \cmidrule(l{3pt}r{3pt}){2-2} \cmidrule(l{3pt}r{3pt}){3-3} \cmidrule(l{3pt}r{3pt}){4-4} \cmidrule(l{3pt}r{3pt}){5-5} \cmidrule(l{3pt}r{3pt}){6-6} \cmidrule(l{3pt}r{3pt}){7-7} \cmidrule(l{3pt}r{3pt}){8-8} \cmidrule(l{3pt}r{3pt}){9-9} \cmidrule(l{3pt}r{3pt}){10-10}
\endfirsthead
\caption[]{Comparing median carbon intensity and horizontal heterogeneity between first and fifth expenditure quintile \textit{(continued)}}\\
\toprule
\endhead

\endfoot
\bottomrule
\insertTableNotes
\endlastfoot
Argentina & 1.50 & 0.69 & 3.54 & 1.52 & 1.58 & 0.73 & 2.18 & 2.33 & 2.16\\
Armenia & 0.78 & 0.45 & 2.51 & 1.56 & 1.07 & 0.58 & 1.71 & 1.61 & 1.85\\
Australia & 0.69 & 0.39 & 0.98 & 0.60 & 0.44 & 0.26 & 1.75 & 1.64 & 1.69\\
Austria & 0.33 & 0.21 & 0.75 & 0.40 & 0.38 & 0.19 & 1.53 & 1.90 & 1.95\\
Bangladesh & 0.32 & 0.32 & 0.33 & 0.48 & 0.15 & 0.22 & 1.00 & 0.69 & 0.69\\
Barbados & 0.37 & 0.43 & 1.59 & 1.00 & 0.65 & 0.51 & 0.86 & 1.59 & 1.29\\
Belgium & 0.46 & 0.36 & 0.81 & 0.53 & 0.38 & 0.27 & 1.29 & 1.52 & 1.42\\
Benin & 0.18 & 0.38 & 1.26 & 1.42 & 0.71 & 0.61 & 0.49 & 0.88 & 1.17\\
Bolivia & 0.35 & 0.34 & 0.73 & 0.50 & 0.30 & 0.25 & 1.04 & 1.48 & 1.19\\
Brazil & 0.51 & 0.36 & 1.01 & 0.59 & 0.40 & 0.29 & 1.40 & 1.72 & 1.38\\
Bulgaria & 0.49 & 0.48 & 0.87 & 0.72 & 0.36 & 0.28 & 1.02 & 1.20 & 1.27\\
Burkina Faso & 0.02 & 0.44 & 1.54 & 1.37 & 0.63 & 0.50 & 0.04 & 1.13 & 1.26\\
Cambodia & 0.41 & 0.37 & 1.14 & 0.87 & 0.57 & 0.43 & 1.13 & 1.31 & 1.33\\
Canada & 0.50 & 0.43 & 1.45 & 0.59 & 0.62 & 0.28 & 1.17 & 2.47 & 2.18\\
Chile & 0.60 & 0.36 & 0.71 & 0.34 & 0.31 & 0.17 & 1.70 & 2.08 & 1.81\\
Colombia & 0.42 & 0.28 & 1.70 & 0.69 & 0.80 & 0.31 & 1.51 & 2.48 & 2.58\\
Costa Rica & 0.15 & 0.16 & 0.54 & 0.50 & 0.26 & 0.28 & 0.97 & 1.08 & 0.96\\
Croatia & 0.26 & 0.34 & 0.77 & 1.02 & 0.34 & 0.41 & 0.78 & 0.75 & 0.83\\
Cyprus & 0.29 & 0.21 & 0.56 & 0.38 & 0.27 & 0.17 & 1.37 & 1.50 & 1.56\\
Czechia & 1.03 & 0.92 & 1.60 & 1.21 & 0.63 & 0.52 & 1.11 & 1.32 & 1.22\\
Côte d’Ivoire & 0.06 & 0.19 & 0.78 & 0.80 & 0.26 & 0.30 & 0.31 & 0.97 & 0.87\\
Denmark & 0.16 & 0.14 & 0.38 & 0.31 & 0.18 & 0.12 & 1.09 & 1.23 & 1.56\\
Dominican Republic & 0.22 & 0.30 & 0.54 & 0.84 & 0.23 & 0.44 & 0.75 & 0.64 & 0.52\\
Ecuador & 0.33 & 0.29 & 1.10 & 1.09 & 0.39 & 0.53 & 1.14 & 1.01 & 0.73\\
Egypt & 0.87 & 0.87 & 0.84 & 1.77 & 0.36 & 0.60 & 1.00 & 0.47 & 0.60\\
El Salvador & 0.19 & 0.18 & 1.43 & 0.85 & 0.84 & 0.34 & 1.05 & 1.68 & 2.44\\
Estonia & 0.55 & 0.41 & 1.25 & 0.72 & 0.56 & 0.37 & 1.36 & 1.74 & 1.53\\
Ethiopia & 0.07 & 0.07 & 0.32 & 0.11 & 0.11 & 0.04 & 0.94 & 3.04 & 2.63\\
Finland & 0.33 & 0.27 & 0.56 & 0.44 & 0.29 & 0.21 & 1.24 & 1.28 & 1.37\\
France & 0.28 & 0.23 & 0.90 & 0.43 & 0.35 & 0.21 & 1.20 & 2.11 & 1.71\\
Georgia & 0.64 & 0.72 & 2.55 & 2.24 & 1.11 & 1.24 & 0.88 & 1.14 & 0.89\\
Germany & 0.53 & 0.42 & 0.89 & 0.60 & 0.43 & 0.27 & 1.26 & 1.48 & 1.58\\
Ghana & 0.07 & 0.16 & 0.50 & 1.16 & 0.11 & 0.24 & 0.42 & 0.44 & 0.47\\
Greece & 0.41 & 0.33 & 0.56 & 0.36 & 0.27 & 0.17 & 1.27 & 1.59 & 1.58\\
Guatemala & 0.06 & 0.35 & 0.39 & 1.12 & 0.12 & 0.54 & 0.16 & 0.35 & 0.22\\
Guinea-Bissau & 0.05 & 0.16 & 0.63 & 0.93 & 0.15 & 0.27 & 0.29 & 0.67 & 0.56\\
Hungary & 0.54 & 0.60 & 1.35 & 1.05 & 0.69 & 0.49 & 0.91 & 1.29 & 1.40\\
India & 1.06 & 1.36 & 1.09 & 1.32 & 0.57 & 0.65 & 0.78 & 0.83 & 0.88\\
Indonesia & 1.33 & 1.33 & 2.58 & 2.13 & 1.33 & 1.05 & 1.00 & 1.21 & 1.27\\
Iraq & 1.05 & 0.62 & 1.87 & 1.29 & 0.84 & 0.60 & 1.70 & 1.45 & 1.39\\
Ireland & 0.46 & 0.32 & 0.96 & 0.52 & 0.41 & 0.23 & 1.44 & 1.86 & 1.77\\
Israel & 0.49 & 0.25 & 0.83 & 0.37 & 0.39 & 0.18 & 1.96 & 2.21 & 2.19\\
Italy & 0.44 & 0.34 & 0.59 & 0.39 & 0.28 & 0.18 & 1.31 & 1.52 & 1.56\\
Jordan & 0.75 & 1.24 & 1.72 & 1.98 & 0.82 & 1.25 & 0.60 & 0.87 & 0.66\\
Kenya & 0.27 & 0.41 & 0.92 & 0.99 & 0.41 & 0.46 & 0.68 & 0.93 & 0.90\\
Latvia & 0.63 & 0.31 & 1.76 & 0.93 & 1.08 & 0.50 & 2.04 & 1.89 & 2.18\\
Liberia & 0.06 & 0.17 & 0.36 & 0.91 & 0.08 & 0.27 & 0.35 & 0.40 & 0.28\\
Lithuania & 0.21 & 0.36 & 1.15 & 1.11 & 0.56 & 0.54 & 0.58 & 1.03 & 1.03\\
Luxembourg & 0.36 & 0.21 & 0.83 & 0.46 & 0.39 & 0.21 & 1.77 & 1.81 & 1.85\\
Malawi & 0.01 & 0.02 & 0.02 & 0.32 & 0.01 & 0.04 & 0.37 & 0.07 & 0.15\\
Maldives & 0.28 & 0.18 & 0.49 & 0.34 & 0.21 & 0.17 & 1.59 & 1.44 & 1.27\\
Mali & 0.02 & 0.26 & 1.03 & 0.78 & 0.54 & 0.45 & 0.06 & 1.31 & 1.20\\
Mexico & 0.70 & 0.85 & 1.64 & 1.45 & 0.74 & 0.72 & 0.82 & 1.13 & 1.03\\
Mongolia & 0.75 & 0.61 & 3.52 & 1.97 & 1.71 & 0.88 & 1.23 & 1.78 & 1.93\\
Morocco & 0.63 & 0.53 & 0.75 & 0.81 & 0.32 & 0.39 & 1.19 & 0.93 & 0.82\\
Mozambique & 0.03 & 0.14 & 1.24 & 1.54 & 0.21 & 0.44 & 0.19 & 0.81 & 0.47\\
Myanmar (Burma) & 0.24 & 0.44 & 0.86 & 1.68 & 0.38 & 0.67 & 0.55 & 0.51 & 0.56\\
Netherlands & 0.40 & 0.31 & 0.36 & 0.27 & 0.18 & 0.14 & 1.28 & 1.31 & 1.27\\
Nicaragua & 0.03 & 0.21 & 0.31 & 0.83 & 0.12 & 0.29 & 0.13 & 0.37 & 0.41\\
Niger & 0.03 & 0.08 & 0.09 & 0.75 & 0.03 & 0.31 & 0.34 & 0.12 & 0.08\\
Nigeria & 0.18 & 0.56 & 0.86 & 1.00 & 0.41 & 0.45 & 0.32 & 0.86 & 0.90\\
Norway & 0.31 & 0.20 & 1.34 & 0.50 & 0.51 & 0.22 & 1.52 & 2.66 & 2.30\\
Pakistan & 0.30 & 0.75 & 0.88 & 1.11 & 0.43 & 0.57 & 0.40 & 0.79 & 0.76\\
Paraguay & 0.23 & 0.28 & 1.01 & 0.65 & 0.49 & 0.28 & 0.82 & 1.55 & 1.73\\
Peru & 0.29 & 0.25 & 0.77 & 0.48 & 0.39 & 0.21 & 1.17 & 1.61 & 1.83\\
Philippines & 0.26 & 0.52 & 0.50 & 0.71 & 0.21 & 0.34 & 0.50 & 0.71 & 0.61\\
Poland & 0.82 & 0.63 & 1.69 & 1.86 & 0.62 & 0.37 & 1.29 & 0.91 & 1.66\\
Portugal & 0.46 & 0.32 & 0.53 & 0.36 & 0.27 & 0.18 & 1.45 & 1.46 & 1.49\\
Romania & 0.49 & 0.43 & 1.31 & 0.77 & 0.54 & 0.37 & 1.14 & 1.71 & 1.49\\
Russia & 1.07 & 1.11 & 2.59 & 1.85 & 1.00 & 0.97 & 0.96 & 1.40 & 1.04\\
Rwanda & 0.00 & 0.02 & 0.04 & 0.68 & 0.00 & 0.03 & 0.17 & 0.06 & 0.10\\
Senegal & 0.07 & 0.26 & 0.48 & 0.91 & 0.14 & 0.26 & 0.27 & 0.53 & 0.53\\
Serbia & 0.46 & 0.46 & 0.54 & 1.81 & 0.25 & 0.25 & 1.01 & 0.30 & 0.99\\
Slovakia & 0.52 & 0.40 & 2.15 & 1.13 & 0.95 & 0.47 & 1.31 & 1.90 & 2.02\\
South Africa & 1.86 & 1.93 & 2.70 & 2.37 & 1.18 & 1.07 & 0.96 & 1.14 & 1.10\\
Spain & 0.40 & 0.37 & 0.67 & 0.48 & 0.33 & 0.22 & 1.08 & 1.41 & 1.47\\
Suriname & 0.11 & 0.09 & 0.36 & 0.27 & 0.15 & 0.09 & 1.26 & 1.31 & 1.64\\
Sweden & 0.15 & 0.14 & 0.57 & 0.32 & 0.26 & 0.14 & 1.06 & 1.79 & 1.92\\
Switzerland & 0.13 & 0.12 & 0.36 & 0.28 & 0.16 & 0.12 & 1.14 & 1.27 & 1.34\\
Taiwan & 1.13 & 1.16 & 0.94 & 1.12 & 0.47 & 0.64 & 0.98 & 0.84 & 0.74\\
Thailand & 1.44 & 1.48 & 2.29 & 2.06 & 1.24 & 1.23 & 0.98 & 1.11 & 1.01\\
Togo & 0.00 & 0.15 & 1.14 & 1.64 & 0.03 & 0.76 & 0.02 & 0.69 & 0.03\\
Turkey & 0.92 & 0.85 & 3.23 & 1.57 & 1.51 & 0.61 & 1.09 & 2.06 & 2.48\\
Uganda & 0.04 & 0.09 & 1.14 & 1.09 & 0.39 & 0.27 & 0.42 & 1.05 & 1.46\\
United Kingdom & 0.51 & 0.34 & 1.81 & 0.74 & 0.86 & 0.32 & 1.48 & 2.44 & 2.67\\
United States & 0.94 & 0.64 & 1.72 & 0.87 & 0.81 & 0.40 & 1.48 & 1.97 & 2.04\\
Uruguay & 0.12 & 0.09 & 0.35 & 0.22 & 0.14 & 0.11 & 1.26 & 1.61 & 1.28\\
Vietnam & 0.54 & 0.44 & 0.49 & 0.39 & 0.23 & 0.19 & 1.23 & 1.24 & 1.18\\*
\end{longtable}
\end{ThreePartTable}
\endgroup{}

\clearpage

% Table C10: Evaluation of BRT
\begingroup\fontsize{8}{10}\selectfont

\begin{ThreePartTable}
\begin{TableNotes}
\item \textit{Note: } 
\item This table shows performance metrics for boosted regression tree models including exclusively household expenditures ('Sparse model') and including all available features ('Rich model'). MAE is the mean absolute error of predictions; RMSE is the root mean squared error of predictions; R\textsuperscript{2} is the squared correlation of prediction errors. Unit of MAE and RMSE is kgCO\textsubscript{2} per USD. We show MAE, RMSE and R\textsuperscript{2} for five-fold cross-validation on the entire dataset. 
\end{TableNotes}
\begin{longtable}[t]{l|r|rrr|rrrl|r|rrr|rrrl|r|rrr|rrrl|r|rrr|rrrl|r|rrr|rrrl|r|rrr|rrrl|r|rrr|rrrl|r|rrr|rrr}
\caption{\label{tab:A8}Evaluation of boosted regression tree models}\\
\toprule
\multicolumn{2}{c}{ } & \multicolumn{3}{c}{Sparse model} & \multicolumn{3}{c}{Rich model} \\
\cmidrule(l{3pt}r{3pt}){3-5} \cmidrule(l{3pt}r{3pt}){6-8}
Country & Mean & MAE & RMSE & R\textsuperscript{2} & MAE & RMSE & R\textsuperscript{2}\\
\midrule
\endfirsthead
\caption[]{Evaluation of boosted regression tree models \textit{(continued)}}\\
\toprule
Country & Mean & MAE & RMSE & R\textsuperscript{2} & MAE & RMSE & R\textsuperscript{2}\\
\midrule
\endhead

\endfoot
\bottomrule
\insertTableNotes
\endlastfoot
Argentina & 1.27 & 0.56 & 0.78 & 0.19 & 0.50 & 0.73 & 0.28\\
Armenia & 0.87 & 0.51 & 0.73 & 0.04 & 0.43 & 0.64 & 0.26\\
Australia & 0.58 & 0.18 & 0.24 & 0.20 & 0.17 & 0.23 & 0.26\\
Austria & 0.31 & 0.14 & 0.19 & 0.05 & 0.12 & 0.17 & 0.21\\
Bangladesh & 0.33 & 0.09 & 0.14 & 0.03 & 0.08 & 0.13 & 0.16\\
Barbados & 0.58 & 0.32 & 0.45 & 0.01 & 0.26 & 0.40 & 0.23\\
Belgium & 0.44 & 0.16 & 0.21 & 0.07 & 0.15 & 0.21 & 0.08\\
Benin & 0.40 & 0.30 & 0.42 & 0.05 & 0.22 & 0.35 & 0.34\\
Bolivia & 0.39 & 0.14 & 0.19 & 0.03 & 0.11 & 0.15 & 0.36\\
Brazil & 0.48 & 0.18 & 0.29 & 0.06 & 0.17 & 0.27 & 0.16\\
Bulgaria & 0.57 & 0.21 & 0.39 & 0.01 & 0.21 & 0.39 & 0.01\\
Burkina Faso & 0.39 & 0.31 & 0.43 & 0.06 & 0.19 & 0.32 & 0.48\\
Cambodia & 0.47 & 0.24 & 0.33 & 0.01 & 0.20 & 0.30 & 0.19\\
Canada & 0.53 & 0.23 & 0.33 & 0.02 & 0.21 & 0.35 & 0.09\\
Chile & 0.50 & 0.12 & 0.18 & 0.18 & 0.12 & 0.17 & 0.24\\
Colombia & 0.50 & 0.27 & 0.43 & 0.04 & 0.24 & 0.41 & 0.16\\
Costa Rica & 0.22 & 0.15 & 0.20 & 0.02 & 0.12 & 0.17 & 0.23\\
Croatia & 0.40 & 0.22 & 0.32 & 0.03 & 0.22 & 0.32 & 0.06\\
Cyprus & 0.27 & 0.10 & 0.14 & 0.05 & 0.10 & 0.14 & 0.08\\
Czechia & 1.11 & 0.34 & 0.51 & 0.03 & 0.33 & 0.49 & 0.10\\
Côte d’Ivoire & 0.22 & 0.17 & 0.25 & 0.05 & 0.12 & 0.20 & 0.40\\
Denmark & 0.18 & 0.08 & 0.11 & 0.01 & 0.08 & 0.11 & 0.05\\
Dominican Republic & 0.30 & 0.16 & 0.22 & 0.05 & 0.11 & 0.17 & 0.41\\
Ecuador & 0.40 & 0.23 & 0.39 & 0.05 & 0.14 & 0.30 & 0.42\\
Egypt & 0.94 & 0.25 & 0.38 & 0.04 & 0.21 & 0.31 & 0.39\\
El Salvador & 0.37 & 0.25 & 0.35 & 0.02 & 0.18 & 0.29 & 0.34\\
Estonia & 0.54 & 0.23 & 0.30 & 0.07 & 0.23 & 0.30 & 0.09\\
Ethiopia & 0.09 & 0.05 & 0.09 & 0.01 & 0.04 & 0.08 & 0.16\\
Finland & 0.33 & 0.13 & 0.16 & 0.00 & 0.12 & 0.16 & 0.09\\
France & 0.32 & 0.17 & 0.28 & 0.01 & 0.14 & 0.24 & 0.25\\
Georgia & 0.96 & 0.62 & 0.83 & 0.02 & 0.49 & 0.72 & 0.27\\
Germany & 0.53 & 0.19 & 0.25 & 0.02 & 0.15 & 0.20 & 0.32\\
Ghana & 0.20 & 0.16 & 0.30 & 0.08 & 0.12 & 0.25 & 0.37\\
Greece & 0.40 & 0.11 & 0.14 & 0.06 & 0.10 & 0.14 & 0.10\\
Guatemala & 0.28 & 0.20 & 0.29 & 0.14 & 0.14 & 0.23 & 0.46\\
Guinea-Bissau & 0.19 & 0.16 & 0.26 & 0.07 & 0.14 & 0.25 & 0.20\\
Hungary & 0.70 & 0.30 & 0.39 & 0.01 & 0.29 & 0.38 & 0.04\\
India & 1.28 & 0.29 & 0.37 & 0.06 & 0.20 & 0.28 & 0.48\\
Indonesia & 1.45 & 0.53 & 0.70 & 0.02 & 0.42 & 0.58 & 0.33\\
Iraq & 0.97 & 0.36 & 0.52 & 0.19 & 0.34 & 0.49 & 0.28\\
Ireland & 0.45 & 0.17 & 0.25 & 0.12 & 0.17 & 0.25 & 0.13\\
Israel & 0.39 & 0.14 & 0.20 & 0.16 & 0.12 & 0.17 & 0.37\\
Italy & 0.41 & 0.12 & 0.16 & 0.06 & 0.11 & 0.16 & 0.10\\
Jordan & 1.10 & 0.49 & 0.61 & 0.03 & 0.29 & 0.40 & 0.58\\
Kenya & 0.40 & 0.23 & 0.36 & 0.04 & 0.21 & 0.34 & 0.12\\
Latvia & 0.52 & 0.37 & 0.47 & 0.07 & 0.32 & 0.44 & 0.20\\
Liberia & 0.19 & 0.14 & 0.24 & 0.09 & 0.14 & 0.24 & 0.11\\
Lithuania & 0.39 & 0.28 & 0.42 & 0.05 & 0.27 & 0.41 & 0.08\\
Luxembourg & 0.32 & 0.15 & 0.20 & 0.15 & 0.15 & 0.20 & 0.16\\
Malawi & 0.03 & 0.03 & 0.12 & 0.10 & 0.03 & 0.11 & 0.21\\
Maldives & 0.26 & 0.10 & 0.13 & 0.07 & 0.09 & 0.12 & 0.20\\
Mali & 0.27 & 0.24 & 0.30 & 0.03 & 0.16 & 0.25 & 0.37\\
Mexico & 0.93 & 0.39 & 0.55 & 0.00 & 0.31 & 0.48 & 0.25\\
Mongolia & 0.98 & 0.64 & 0.95 & 0.04 & 0.57 & 0.87 & 0.21\\
Morocco & 0.62 & 0.17 & 0.25 & 0.02 & 0.17 & 0.24 & 0.07\\
Mozambique & 0.22 & 0.24 & 0.42 & 0.12 & 0.23 & 0.42 & 0.13\\
Myanmar (Burma) & 0.45 & 0.27 & 0.39 & 0.06 & 0.26 & 0.38 & 0.12\\
Netherlands & 0.36 & 0.08 & 0.10 & 0.08 & 0.08 & 0.10 & 0.14\\
Nicaragua & 0.19 & 0.13 & 0.20 & 0.12 & 0.09 & 0.15 & 0.49\\
Niger & 0.09 & 0.09 & 0.18 & 0.18 & 0.06 & 0.14 & 0.49\\
Nigeria & 0.46 & 0.24 & 0.33 & 0.06 & 0.19 & 0.27 & 0.36\\
Norway & 0.33 & 0.19 & 0.30 & 0.05 & 0.16 & 0.26 & 0.26\\
Pakistan & 0.60 & 0.26 & 0.34 & 0.12 & 0.23 & 0.31 & 0.27\\
Paraguay & 0.35 & 0.19 & 0.29 & 0.00 & 0.16 & 0.25 & 0.20\\
Peru & 0.32 & 0.13 & 0.19 & 0.04 & 0.10 & 0.16 & 0.31\\
Philippines & 0.45 & 0.15 & 0.20 & 0.17 & 0.12 & 0.17 & 0.44\\
Poland & 1.09 & 0.61 & 1.31 & 0.01 & 0.61 & 1.29 & 0.03\\
Portugal & 0.40 & 0.11 & 0.15 & 0.05 & 0.10 & 0.14 & 0.15\\
Romania & 0.61 & 0.31 & 0.56 & 0.04 & 0.24 & 0.50 & 0.23\\
Russia & 1.28 & 0.53 & 0.81 & 0.02 & 0.41 & 0.70 & 0.27\\
Rwanda & 0.04 & 0.05 & 0.13 & 0.23 & 0.03 & 0.11 & 0.49\\
Senegal & 0.23 & 0.14 & 0.22 & 0.07 & 0.10 & 0.18 & 0.36\\
Serbia & 0.64 & 0.39 & 1.10 & 0.01 & 0.38 & 1.08 & 0.03\\
Slovakia & 0.71 & 0.45 & 0.72 & 0.03 & 0.41 & 0.70 & 0.11\\
South Africa & 2.03 & 0.62 & 0.87 & 0.02 & 0.50 & 0.71 & 0.34\\
Spain & 0.42 & 0.14 & 0.18 & 0.00 & 0.13 & 0.18 & 0.09\\
Suriname & 0.13 & 0.07 & 0.11 & 0.01 & 0.07 & 0.11 & 0.01\\
Sweden & 0.19 & 0.13 & 0.24 & 0.05 & 0.12 & 0.24 & 0.03\\
Switzerland & 0.15 & 0.08 & 0.11 & 0.01 & 0.07 & 0.10 & 0.18\\
Taiwan & 1.17 & 0.25 & 0.31 & 0.02 & 0.20 & 0.25 & 0.37\\
Thailand & 1.56 & 0.52 & 0.66 & 0.07 & 0.42 & 0.56 & 0.34\\
Togo & 0.26 & 0.31 & 0.47 & 0.06 & 0.19 & 0.37 & 0.41\\
Turkey & 1.25 & 0.62 & 0.87 & 0.04 & 0.54 & 0.80 & 0.20\\
Uganda & 0.20 & 0.20 & 0.38 & 0.06 & 0.14 & 0.33 & 0.30\\
United Kingdom & 0.55 & 0.28 & 0.42 & 0.09 & 0.25 & 0.40 & 0.21\\
United States & 0.86 & 0.29 & 0.41 & 0.09 & 0.28 & 0.39 & 0.14\\
Uruguay & 0.13 & 0.07 & 0.10 & 0.01 & 0.05 & 0.08 & 0.34\\
Vietnam & 0.51 & 0.10 & 0.14 & 0.10 & 0.09 & 0.12 & 0.29\\*
\end{longtable}
\end{ThreePartTable}
\endgroup{}

\clearpage

% Table C11: Feature importance across countries by cluster
\input{2_Tables/Table_Countries_SHAP_Summary_Corrected}
\clearpage

% Table C12: Feature importance across countries by cluster (uncorrected)
\input{2_Tables/Table_Countries_SHAP_Summary_Uncorrected}
\clearpage

% Table C13: Feature importance across countries by cluster (corrected and imputed)
\input{2_Tables/Table_Countries_SHAP_Summary_Corrected_Imputed}
\clearpage

% Table C14: Electricity generation
\input{2_Tables/Table_Electricity_Generation}
\clearpage

% Table C15: Summary of clusters
\begin{table}[H]
\centering
\caption{\label{tab:A9}Average feature importance across country clusters}
\centering
\resizebox{\ifdim\width>\linewidth\linewidth\else\width\fi}{!}{
\begin{threeparttable}
\begin{tabular}[t]{lrrrrrrrrrrrrr}
\toprule
\rotatebox{90}{Cluster} & \rotatebox{90}{Number} & \rotatebox{90}{Average silhouette width} & \rotatebox{90}{Vertical distribution} & \rotatebox{90}{HH expenditures} & \rotatebox{90}{Sociodemographic} & \rotatebox{90}{Spatial} & \rotatebox{90}{Electricity access} & \rotatebox{90}{Cooking fuel} & \rotatebox{90}{Heating fuel} & \rotatebox{90}{Lighting fuel} & \rotatebox{90}{Car own.} & \rotatebox{90}{Motorcycle own.} & \rotatebox{90}{Appliance own.}\\
\midrule
A & 38 & 0.42 & 0.10 & 0.03 & 0.02 & 0.03 & 0.00 & 0.00 & 0.00 & 0.00 & 0.01 & 0.01 & 0.00\\
B & 12 & 0.10 & -0.46 & 0.05 & 0.04 & 0.04 & 0.00 & 0.06 & 0.00 & 0.00 & 0.06 & 0.05 & 0.03\\
C & 8 & 0.12 & -0.09 & 0.07 & 0.04 & 0.05 & 0.00 & 0.01 & 0.00 & 0.00 & 0.17 & 0.03 & 0.02\\
D & 7 & 0.25 & 1.60 & 0.11 & 0.07 & 0.03 & 0.00 & 0.00 & 0.00 & 0.00 & 0.02 & 0.00 & 0.01\\
E & 6 & 0.39 & 0.64 & 0.06 & 0.02 & 0.03 & 0.00 & 0.01 & 0.09 & 0.00 & 0.03 & 0.00 & 0.01\\
F & 6 & 0.38 & -1.62 & 0.04 & 0.06 & 0.04 & 0.01 & 0.01 & 0.00 & 0.01 & 0.04 & 0.19 & 0.01\\
G & 5 & 0.28 & 0.73 & 0.06 & 0.04 & 0.12 & 0.00 & 0.00 & 0.00 & 0.00 & 0.01 & 0.00 & 0.01\\
H & 2 & 0.58 & -1.44 & 0.06 & 0.04 & 0.06 & 0.00 & 0.03 & 0.00 & 0.09 & 0.05 & 0.05 & 0.01\\
I & 2 & 0.52 & -1.44 & 0.02 & 0.04 & 0.05 & 0.04 & 0.07 & 0.00 & 0.00 & 0.04 & 0.07 & 0.04\\
J & 2 & 0.08 & -0.28 & 0.07 & 0.03 & 0.07 & 0.01 & 0.00 & 0.00 & 0.00 & 0.02 & 0.04 & 0.13\\
\bottomrule
\end{tabular}
\begin{tablenotes}
\item \textit{Note: } 
\item This table shows the average importance of features in percent (based on absolute average SHAP-values per feature) across all countries from each cluster A to J. We adjust feature importance for model accuracy. Column 'Vertical distribution' shows average values. Column 'number' refers to the number of countries assigned to this cluster.
\end{tablenotes}
\end{threeparttable}}
\end{table}

\clearpage

% Table C16: Comparing vertical and horizontal distribution coefficients
\input{2_Tables/Table_Vertical_Horizontal_Policy_Simulation}
\clearpage

\begin{refcontext}[sorting=nyt]
\printbibliography[heading=subbibliography, title ={References (Appendix)}]
\end{refcontext}
\end{refsection}

\clearpage

\section{Supplementary information}

\subsection{Data availability} \label{data_availability} Data from household budget surveys are available from statistical agencies subject to permission and possible allowances. See also Table \ref{tab:datasets}. Data from GTAP are available through \href{https://www.gtap.agecon.purdue.edu/}{GTAP}, subject to academic subscription.
Descriptive analyses for carbon pricing reforms and stylized compensation policies can be accessed and customized through \href{https://www.cpic-global.net/}{a separate webtool}.

\subsection{Code availability} \label{code}
We distribute all code written for cleaning and harmonizing household data, modeling carbon intensity of consumption and analysis through \href{https://github.com/lmissbach/Carbon-Intensity-2023}{GitHub}. This repository also contains matching tables for all countries.


\end{document}

